%!TEX TS-program = xelatex
%!TEX TS-options = -shell-escape
%!TEX root = ../LieGrp_SS16/liegruppen.tex
\RequirePackage{fix-cm} 
\documentclass[a4paper, twoside, headsepline, index=totoc,toc=listof,toc=bibliography,toc=index, fontsize=10pt, cleardoublepage=empty, headinclude, DIV=12, BCOR=5mm, titlepage,draft]{scrreprt}
%!TEX root = ../AnaTopGeo_SS14/ana_top_geo.tex
\usepackage{scrtime} % KOMA, Uhrzeit ermoeglicht

%--Pakete zum "Programmieren"
% ======================================================================================
\usepackage{etoolbox}
\usepackage{letltxmacro}
\usepackage{ifthen}
% ======================================================================================

%--Farbdefinitionen und Grafiken (muss vor tikz geladen werden)
% ======================================================================================
\usepackage[usenames, table, x11names]{xcolor}
\definecolor{dark_gray}{gray}{0.45}
\definecolor{light_gray}{gray}{0.6}
\definecolor{fb10_blue}{cmyk}{0.8,0.4,0.13,0.07}
\usepackage[final]{graphicx}
\usepackage{adjustbox}
\newcommand{\cfbox}[2]{% coloured frame box
	\ifmmode
	\mathchoice{\adjustbox{cfbox=#1}{$\displaystyle#2$}}{\adjustbox{cfbox=#1}{$\textstyle#2$}}{\adjustbox{cfbox=#1}{$\scriptstyle#2$}}{\adjustbox{cfbox=#1}{$\scriptscriptstyle#2$}}
	\else
	\adjustbox{cfbox=#1}{#2}
	\fi
}
% ======================================================================================

%--Zum Zeichnen/ TikZ-Kram (vor polyglossia bzw. babel geladen werden)
% ======================================================================================
\usepackage{tikz}
\usepackage{tikz-cd}
\usetikzlibrary{external}
\tikzset{>=latex}
\usetikzlibrary{%
	shapes,
	arrows.meta,
	intersections,
	calc,
	3d,
	decorations.pathreplacing,decorations.markings,decorations.pathmorphing,
	angles,
	quotes,
}
\tikzexternalize[prefix=tikz/,up to date check=diff]
\pgfkeys{/pgf/images/include external/.code=\includegraphics{#1}}
\tikzset{external/system call={lualatex \tikzexternalcheckshellescape -halt-on-error -interaction=batchmode --shell-escape -jobname "\image" "\texsource"}}
\AtBeginEnvironment{tikzcd}{\tikzexternaldisable} % tikzexternalize fuer tikzcd deaktivieren, da inkompatibel
\AtEndEnvironment{tikzcd}{\tikzexternalenable}
\tikzset{% um Inkompatibilitaeten von quotes und polyglossia bzw. babel zu vermeiden
  every picture/.append style={
    execute at begin picture={\shorthandoff{"}},
    execute at end picture={\shorthandon{"}}
  }
}
\usepackage{pgfplots}
\usepgfplotslibrary{colormaps}
\newcommand*\circled[1]{\tikzexternaldisable\tikz[baseline=(char.base)]{\node[shape=circle,draw,inner sep=2pt] (char) {#1};}\tikzexternalenable}
% ======================================================================================



%-- Mathepakete etc.
% ======================================================================================
\usepackage[T1]{fontenc}
\renewcommand{\rmdefault}{zpltlf}
\usepackage{mathtools} % beinhaltet amsmath
\mathtoolsset{showonlyrefs,centercolon,showmanualtags}
\newtagform{brackets}[\textbf]{[}{]}
\usetagform{brackets}
\usepackage{fix-cm}
\usepackage[bbgreekl]{mathbbol}
\usepackage{amssymb,marvosym} 
\usepackage{nicefrac} % schräge Brüche
\usepackage{faktor}
\newcommand{\Faktor}[1]{\faktor[\textstyle]{#1}}
\usepackage{xfrac}
\usepackage{cancel}
\usepackage{mathdots} % Verbesserung von Punkten wie zB \ldots
\usepackage[bb=px]{mathalfa} % \mathbb als px font
\usepackage{centernot}
\usepackage{stackrel}
\DeclareSymbolFont{bbold}{U}{bbold}{m}{n}
\DeclareSymbolFontAlphabet{\mathbbold}{bbold}
\newcommand{\ind}{\mathbbold{1}} % charakteristische-Funktion-Eins
\def\mathul#1#2{\color{#1}\underline{{\color{black}#2}}\color{black}} %farbiges Untersteichen im Mathe-Modus
\renewcommand{\le}{\leqslant}
\renewcommand{\ge}{\geqslant}
% ======================================================================================


%-- Von xfrac erzeuge font warnings ignorieren
% ======================================================================================
\usepackage{silence}
\WarningFilter{latexfont}{Size substitutions with differences}
\WarningFilter{latexfont}{Font shape `U/bbold/m/n' in size}
% ======================================================================================


%-- Typographie/Polyglossia
% ======================================================================================
\usepackage[euler-digits]{eulervm} % vor fontspec laden!
\usepackage[no-math]{fontspec}
\usepackage{polyglossia} % moderner babel-ersatz
\setmainlanguage[spelling=new,babelshorthands=true]{german}
\shorthandoff{"}
\setotherlanguage{english}
\defaultfontfeatures{Mapping=tex-text, WordSpace={1.2}, Ligatures={Required,Common,Contextual},Extension=.otf} %


\setmainfont{TeXGyrePagellaX}[UprightFont=*-Regular,BoldFont=*-Bold,ItalicFont=*-Italic,BoldItalicFont=*-BoldItalic,ItalicFeatures={Style=Historic},Ligatures={Required,Common,Contextual,Historic}]
\setsansfont{texgyreadventor}[Scale=MatchUppercase, UprightFont=*-regular, BoldFont=*-bold, ItalicFont=*-italic, BoldItalicFont=*-bolditalic]
\setmonofont{SourceCodePro}[Scale=0.9,UprightFont=*-Regular, BoldFont=*-Semibold, ItalicFont=*-Light]
\usepackage{xltxtra}
\usepackage{fontawesome}
\usepackage[final]{microtype}
\usepackage[draft=false]{scrlayer-scrpage} 
\flushbottom
% ======================================================================================


%-- Aufzählungen
% ======================================================================================
\usepackage[shortlabels,inline]{enumitem}
\setlist[itemize,1]{label=\faCaretRight}
\setlist[enumerate]{font=\bfseries}
\setlist[description]{font=\normalfont\bfseries}
\usepackage{multicol}
% ======================================================================================


%-- Floats/Figures/Tabellen
% ======================================================================================
\usepackage{wrapfig}
\usepackage{float}
\usepackage[margin=10pt, font=small, labelfont={sf, bf}, format=plain, indention=1em]{caption}
\captionsetup[wrapfigure]{name=Abb. }
\usepackage{booktabs}
% ======================================================================================


%-- korrekte Anführungszeichen und Zitierbefehle
% ======================================================================================
\usepackage[autostyle,german=quotes,english=british]{csquotes}
% ======================================================================================


%--Indexverarbeitung
% ======================================================================================
\usepackage{makeidx}
\newcommand{\bet}[1]{\textbf{\emph{#1}}}
\newcommand{\Index}[1]{\bet{#1}\index{#1}}
\makeindex
\setindexpreamble{{\noindent\sffamily\small Die \emph{Seitenzahlen} sind mit Hyperlinks versehen und somit anklickbar} \par \bigskip}
\renewcommand{\indexpagestyle}{scrheadings}
% ======================================================================================


%-- Marginnotes/Todonotes/Footnotes
% ======================================================================================
\deffootnote[1.5em]{1.5em}{1.5em}{\textsuperscript{\thefootnotemark}\ }
\usepackage[fulladjust]{marginnote}
\renewcommand*{\marginfont}{\itshape\footnotesize}
\usepackage[textsize=small]{todonotes}
\usepackage{ragged2e}
\renewcommand*{\raggedleftmarginnote}{\RaggedLeft}
\renewcommand*{\raggedrightmarginnote}{\RaggedRight}
\LetLtxMacro{\oldtodo}{\todo}
\renewcommand{\todo}[2][]{\tikzexternaldisable\oldtodo[#1]{#2}\tikzexternalenable}
\LetLtxMacro{\oldmissingfigure}{\missingfigure}
\renewcommand{\missingfigure}[2][]{\tikzexternaldisable\oldmissingfigure[{#1}]{#2}\tikzexternalenable}
% ======================================================================================


% -- BibLaTeX
% ======================================================================================
\usepackage[%
	backend=biber,
	sortlocale=auto,
	natbib,
	hyperref,
	backref,
	style=alphabetic
	]%
{biblatex}
\renewcommand*{\mkbibnamelast}[1]{%
  \ifmknamesc{\textsc{#1}}{#1}}
\renewcommand*{\mkbibnameprefix}[1]{%
  \ifboolexpr{ test {\ifmknamesc} and test {\ifuseprefix} }
    {\textsc{#1}}
    {#1}}
\def\ifmknamesc{%
  \ifboolexpr{ test {\ifcurrentname{labelname}}
               or test {\ifcurrentname{author}}
               or ( test {\ifnameundef{author}} and test {\ifcurrentname{editor}} ) }}
\addbibresource{../!config/quellen.bib}
% ======================================================================================

%--Konfiguration von Hyperref und Cleveref
% ======================================================================================
\usepackage[hidelinks, pdfpagelabels,  bookmarksopen=true, bookmarksnumbered=true, linkcolor=black, urlcolor=SkyBlue2, plainpages=false,pagebackref, citecolor=black, hypertexnames=true, pdfauthor={Jannes Bantje}, pdfborderstyle={/S/U}, linkbordercolor=SkyBlue2, colorlinks=false,final,backref=false]{hyperref}
\usepackage[nameinlink,noabbrev]{cleveref}
\newcommand{\appendLink}[1]{#1\,\faExternalLink}
\newcommand{\hrefsym}[2]{\href{#1}{\texttt{\appendLink{#2}}}}
\newcommand{\hrefsymX}[2]{\href{#1}{\appendLink{#2}}}
\newcommand{\hrefsymmail}[2]{\href{#1}{\texttt{\faEnvelopeO\,#2}}}
\renewcommand{\url}[1]{\hrefsym{#1}{\nolinkurl{#1}}}
% ======================================================================================


% -- QR-Codes (hinter hyperref laden!)
% ======================================================================================
\usepackage{qrcode}
% ======================================================================================

%--Römische Zahlen
% ======================================================================================
\newcommand{\RM}[1]{\MakeUppercase{\romannumeral #1{}}}
% ======================================================================================

%-- Definition von diversen Mathe-Befehlen
% ======================================================================================
%!TEX root = mitschrift_main.tex

% -- Zum Finetuning von Befehlen
% ======================================================================================
\makeatletter
\newcommand{\raisemath}[1]{\mathpalette{\raisem@th{#1}}}
\newcommand{\raisem@th}[3]{\raisebox{#1}{$#2#3$}}
\makeatother
\makeatletter
\newcommand{\killDescendersM}[1]{\mathpalette{\killD@scendersM{#1}}}
\newcommand{\killD@scendersM}[2]{\raisebox{0pt}[\height][0pt]{$#2#1$}}
\makeatother
\DeclareRobustCommand{\minwidthbox}[2]{%
  \ifmmode
    \expandafter\mathmakebox
  \else
    \expandafter\makebox
  \fi
  [\ifdim#2<\width\width\else#2\fi]{#1}%
}
% ======================================================================================


%-- Klammerbefehle
% ======================================================================================
\DeclarePairedDelimiter{\abs}{\lvert}{\rvert}
\DeclarePairedDelimiter{\floor}{\lfloor}{\rfloor}
\DeclarePairedDelimiter{\ceil}{\lceil}{\rceil}
\DeclarePairedDelimiter\norm{\Vert}{\Vert}
\DeclarePairedDelimiter\enbrace{(}{)}
\DeclarePairedDelimiter\benbrace{[}{]}
\DeclarePairedDelimiter\bbenbrace{[\![}{]\!]}
\DeclarePairedDelimiter\lenbrace{<}{>}
\DeclarePairedDelimiter\angbrace{\langle}{\rangle}
\newcommand{\ssbrace}[1]{{\scriptscriptstyle\enbrace{#1}}}
\newcommand{\ssbbrace}[1]{{\scriptscriptstyle\benbrace{#1}}}
% ======================================================================================

%-- Mengen
% ======================================================================================
\newcommand\SetSymbol[1][]{\nonscript\:#1\vert\allowbreak\nonscript\:\mathopen{}}
\providecommand\given{} % to make it exist
\DeclarePairedDelimiterX\set[1]\{\}{\renewcommand\given{\SetSymbol[\delimsize]}#1}
% ======================================================================================

%-- Skalarprodukt (3 Varianten) 
% ======================================================================================
\DeclarePairedDelimiterX\sprod[2]{\langle}{\rangle}{#1\,\delimsize\vert\,#2}
\DeclarePairedDelimiterX\skal[2]{\langle}{\rangle}{#1\,,\,#2}
\makeatletter
\DeclareFontFamily{OMX}{MnSymbolE}{}
\DeclareSymbolFont{MnLargeSymbols}{OMX}{MnSymbolE}{m}{n}
\SetSymbolFont{MnLargeSymbols}{bold}{OMX}{MnSymbolE}{b}{n}
\DeclareFontShape{OMX}{MnSymbolE}{m}{n}{
    <-6>  MnSymbolE5
   <6-7>  MnSymbolE6
   <7-8>  MnSymbolE7
   <8-9>  MnSymbolE8
   <9-10> MnSymbolE9
  <10-12> MnSymbolE10
  <12->   MnSymbolE12
}{}
\DeclareFontShape{OMX}{MnSymbolE}{b}{n}{
    <-6>  MnSymbolE-Bold5
   <6-7>  MnSymbolE-Bold6
   <7-8>  MnSymbolE-Bold7
   <8-9>  MnSymbolE-Bold8
   <9-10> MnSymbolE-Bold9
  <10-12> MnSymbolE-Bold10
  <12->   MnSymbolE-Bold12
}{}
\let\llangle\@undefined
\let\rrangle\@undefined
\DeclareMathDelimiter{\llangle}{\mathopen}%
                     {MnLargeSymbols}{'164}{MnLargeSymbols}{'164}
\DeclareMathDelimiter{\rrangle}{\mathclose}%
                     {MnLargeSymbols}{'171}{MnLargeSymbols}{'171}
\makeatother
\DeclarePairedDelimiterX\sskal[2]{\llangle}{\rrangle}{#1\,,\,#2}
% ======================================================================================

%-- Abbildungsdefinition
% ======================================================================================
\newcommand{\mapdef}[5]{%
	\[
		\begin{array}{rcl}
			\textstyle #1 &\xrightarrow{\minwidthbox{#5}{2em}} & \textstyle #2 \\[0.5ex]
			\textstyle #3 &\xmapsto{\minwidthbox{\mbox{ }}{2em}} & \textstyle #4
		\end{array}
	\]
}
% ======================================================================================

%-- modifiziertes Stackrel 
% ======================================================================================
\newcommand{\StackText}[2]{\stackrel{\mbox{\scriptsize #1}}{#2}}
\newcommand{\StackTextClap}[2]{\stackrel{\mathclap{\mbox{\scriptsize #1}}}{#2}}
% ======================================================================================

%-- Blitz
% ======================================================================================
\newcommand{\light}{\text{\raisebox{-.3ex}{\Large\Lightning}}}
% ======================================================================================


%-- Underbrace u.Ä. als Befehl in LaTeX-Syntax (und ohne Spacingprobleme mit nachfolgenden Operatoren...)
% ======================================================================================
\newcommand{\Underbrace}[2]{{\underbrace{#1}_{#2}}}
\newcommand{\Underbracket}[2]{{\underbracket[0.7pt][2pt]{#1}_{#2}}}
\newcommand{\Overbracket}[2]{{\overbracket[0.7pt][2pt]{#1}^{#2}}}
% ======================================================================================


%-- Deklaration weiterer Operatoren (allgemein)
% ======================================================================================
\DeclareMathOperator{\re}{Re} % Realteil
\let\Re\relax
\DeclareMathOperator{\Re}{Re} % Realteil
\DeclareMathOperator{\im}{im} % Bild
\let\Im\relax
\DeclareMathOperator{\Im}{Im} % Bild
\DeclareMathOperator{\id}{id} % identische Abbildung
\DeclareMathOperator{\conj}{conj} % Konjugation
\DeclareMathOperator{\sgn}{sgn} % Signum
\DeclareMathOperator{\End}{End} % Endomorphismen
\DeclareMathOperator{\Hom}{Hom} % Homomorphismen
\DeclareMathOperator{\Iso}{Iso} % Isomorphismen
\DeclareMathOperator{\Aut}{Aut} % Automorphismen
\DeclareMathOperator{\Span}{span} % Span
\DeclareMathOperator{\coker}{coker} % Kokern
\DeclareMathOperator{\Tr}{Tr} % Spur,Trace
\DeclareMathOperator{\pr}{pr} % Projektion
\DeclareMathOperator{\diag}{diag} % Diagonalmatrix
\DeclareMathOperator{\Rg}{Rg} % Rang
\DeclareMathOperator{\const}{const} % konstante Abbildung
\DeclareMathOperator{\Spur}{Spur} % Spur
\DeclareMathOperator{\Arg}{Arg} % Argument
\DeclareMathOperator{\dist}{dist} % Distanz
\DeclareMathOperator{\supp}{supp} % Träger
\DeclareMathOperator{\Char}{char} % Charakteristik
% ======================================================================================


%-- Deklaration weiterer Operatoren (Differentiale etc.)
% ======================================================================================
\DeclareMathOperator{\grad}{grad} % Gradient
\DeclareMathOperator{\dive}{div} % Gradient
\DeclareMathOperator{\rot}{rot} % Rotation
\newcommand{\D}{\ensuremath{\mathrm{D}\mkern-1.0mu}} % Differential
\newcommand{\mathd}{\ensuremath{\mathrm{d}\mkern-1.0mu}} % äußere Ableitung
\newcommand{\Tmap}{\ensuremath{\mathrm{T}\mkern-0.85mu}} % Tangentialraum
\let\Tang\Tmap
\DeclareMathOperator{\Diff}{Diff}
\newcommand{\diff}[2]{\ensuremath{\frac{{\partial #1}}{{\partial #2}} }}
\newcommand{\diffd}[2]{\ensuremath{\frac{\mathd #1}{\mathd #2} }}
% ======================================================================================


%-- Deklaration weiterer Operatoren (Topologie)
% ======================================================================================
\newcommand*\interior[1]{\overset{\smash{\raisebox{-0.18ex}{$\scriptstyle\circ$}}}{#1}}
\newcommand{\sing}{{\raisemath{1.1pt}{\scriptscriptstyle\mathrm{sing}}}}
\newcommand{\pt}{\mathrm{pt}}
\DeclareMathOperator{\Zyl}{Zyl}
\DeclareMathOperator{\Tel}{Tel}
\newcommand{\op}{\mathrm{op}}
\DeclareMathOperator{\Sp}{Sp}
\DeclareMathOperator{\Keg}{Keg}
\newcommand{\slashedi}{i\hspace{-3.5pt}/}
\newcommand{\cupp}{\smallsmile}
\newcommand{\capp}{\smallfrown}
\DeclareMathOperator*{\colim}{colim}
\DeclareMathOperator{\PD}{PD}
\newcommand{\lf}{\mathrm{lf}}
\DeclareMathOperator{\sig}{sig}
\DeclareMathOperator{\Tor}{Tor}
\DeclareMathOperator{\Ext}{Ext}
\DeclareMathOperator{\AW}{AW}
\DeclareMathOperator{\Proj}{Proj}
\DeclareMathOperator{\Gr}{Gr}
\DeclareMathOperator{\res}{res}
\DeclareMathOperator{\Spec}{Spec}
\DeclareMathOperator{\co}{co}
\DeclareMathOperator{\ch}{ch}
\DeclareMathOperator{\wOp}{w}
\DeclareMathOperator{\Ar}{Ar}
\newcommand{\actson}{\mathrel{\curvearrowright}}
\let\acts\actson
\let\action\actson
\DeclareMathSymbol{\bbDelta}{\mathord}{bbold}{"01}
\newcommand{\DDelta}{\bbDelta}
\DeclareMathOperator{\Star}{Star}
\DeclareMathOperator{\Link}{Link}
\DeclareMathOperator{\EPK}{EPK}
\DeclareMathOperator{\Vol}{Vol}
\newcommand{\cell}{{\raisemath{1.1pt}{\scriptscriptstyle\mathrm{cell}}}}
\DeclarePairedDelimiter{\homologieklasse}{\llbracket}{\rrbracket}
\newcommand{\rand}[1]{\ensuremath{\partial^{\scriptscriptstyle #1}}}
\DeclareMathOperator{\ab}{ab}
% ======================================================================================


%-- Deklaration von Operatoren (Liegruppen)
% ======================================================================================
\DeclareMathOperator{\GL}{GL}
\DeclareMathOperator{\SO}{SO}
\DeclareMathOperator{\Ad}{Ad}
\DeclareMathOperator{\ad}{ad}
\DeclareMathOperator{\On}{O}
\DeclareMathOperator{\Un}{U}
\DeclareMathOperator{\SU}{SU}
\DeclareMathOperator{\Mat}{Mat}
\DeclareRobustCommand{\Der}{\mathop{\mathfrak{der}}}
\DeclareMathOperator{\SL}{SL}
\DeclareMathOperator{\Graph}{Graph}
\DeclareMathOperator{\Int}{Int}
\DeclareRobustCommand{\intAlg}{\mathop{\mathfrak{int}}}
\DeclareMathOperator{\aut}{aut}
\DeclareMathOperator{\Rad}{Rad}
\DeclareMathOperator{\Nil}{Nil}
\DeclareMathOperator{\rad}{rad}
\DeclareMathOperator{\nil}{nil}
\DeclareMathOperator{\Ric}{Ric}
\DeclareMathOperator{\ric}{ric}
\newcommand{\bi}{\mathrm{bi}}
\DeclareMathOperator{\Isom}{Isom}
\DeclareMathOperator{\Sym}{Sym}
\newcommand{\opL}{\ensuremath{\mathrm{L}\mkern-0.6mu}}
% ======================================================================================

%-- Deklaration von Operatoren (Funktionalanalysis)
% ======================================================================================
\DeclareMathOperator{\tr}{tr}
\newcommand{\w}{\mkern1mu\mathrm{w}}
\newcommand{\sa}{\mathrm{sa}}
\newcommand{\vb}{\mathrm{v\mkern-2.5mu.b\mkern-1.5mu.}} % vollständig beschränkt
\newcommand{\so}{\mathrm{\mkern.3mu s\mkern-1.4mu.\mkern-.6mu o\mkern-1.7mu.}} % \newcommand{\so}{\mathrm{s.o.}}
\newcommand{\solim}{\so\text{-}\mkern-0.8mu\lim}
\newcommand{\wo}{\mathrm{w\mkern-3mu.\mkern-.4mu o\mkern-1.7mu.}}
\newcommand{\Top}[1]{\mathcal{T}_{\mkern-2.3mu #1}}
\newcommand{\weakT}[1]{\ensuremath{\mathcal{T}_{#1}^{\mkern+1.0mu\text{\raisebox{0.4ex}{$\mathrm{w}$}}}}}
\newcommand{\weakTstar}[1]{\ensuremath{\mathcal{T}_{#1}^{\mkern+1.0mu\text{\raisebox{0.4ex}{$\mathrm{w}$}}^*}}}
\newcommand{\TWeakStar}{\Top{\w^*}}
\newcommand{\TWeakOp}{\Top{\wo}}
\newcommand{\Tso}{\Top{\so}}
\newcommand{\finSub}{\subset\mkern-0.7mu \subset}
\DeclareMathOperator{\Inv}{Inv}
\newcommand{\simm}{{\hspace{-1.6pt}\raisemath{0.5pt}{\sim}}}
\newcommand{\plus}{{\hspace{-1.6pt}+}}
\DeclareMathOperator{\ev}{ev}
\DeclareMathOperator{\Alg}{Alg}
\DeclareMathOperator{\her}{her}
\newcommand{\subher}{\subset_{\her}}
\newcommand{\grenzw}[1]{\xrightarrow{\minwidthbox{#1}{1.4em}}}
\newcommand{\grenzwl}[1]{\xleftarrow{\minwidthbox{#1}{1.4em}}}
\newcommand{\grenzwIn}[1]{\grenzw{\raisemath{-2pt}{#1}}}
\newcommand{\MyTo}[1]{\tikzexternaldisable\mathbin{\tikz[baseline] \draw[-to,line width=.4pt] (0ex,0.94ex) -- (#1,0.94ex);}\tikzexternalenable}
\newcommand{\dlim}{%
    \mathchoice
      {\lim\limits_{\MyTo{4.2ex}}}% \displaystyle
      {\lim\limits_{\MyTo{2.8ex}}}% \textstyle
      {\lim\limits_{\MyTo{2.3ex}}}% \scriptstyle
      {\lim\limits_{\MyTo{2.3ex}}}% \scriptscriptstyle
}
\newcommand{\Dlim}{\killDescendersM{\dlim}}
\DeclareMathOperator{\sep}{sep}
\DeclareMathOperator{\diam}{diam}
\DeclareMathOperator{\conv}{conv}
\DeclareMathOperator{\Prim}{Prim}
\DeclareMathOperator{\hull}{hull}
\DeclareMathOperator{\red}{red}
\DeclarePairedDelimiterX\bra[1]{\langle}{\rvert}{#1\,}
\DeclarePairedDelimiterX\ket[1]{\lvert}{\rangle}{\,#1}
\DeclarePairedDelimiterX\bracket[2]{\langle}{\rangle}{#1\,\delimsize\vert\,#2}
\newcommand{\tensormax}{\mathbin{\otimes_{\max}}}
\newcommand{\tensormin}{\mathbin{\otimes_{\min}}}
\DeclareMathOperator{\Ped}{Ped}
\newcommand{\alg}{\mathrm{alg}}
\DeclareMathOperator{\CPC}{CPC}
\DeclareMathOperator{\CP}{CP}
\DeclareMathOperator{\UPC}{UPC}
\newcommand{\DeltaOp}{\mathbin{\Delta}}
\newcommand{\kernedP}{\mathcal{P}\mkern-2mu}
\newcommand{\Pinfty}{\kernedP_{\infty}}
\DeclareMathOperator{\Groth}{Groth}
\DeclareMathOperator{\rk}{rk}
\newcommand{\MvN}{\mathrm{MvN}}
% ======================================================================================

%-- Kategorien
% ======================================================================================
\DeclareMathOperator{\Mor}{Mor}
\DeclareMathOperator{\Obj}{Obj}
\newcommand{\TOP}{\textsc{Top}}
\newcommand{\HTOP}{\textsc{HTop}}
\newcommand{\VR}{\textsc{VR}}
\newcommand{\MOD}{\textsc{Mod}}
\newcommand{\Mod}[1]{#1\text{-}\MOD}
\newcommand{\MONOIDE}{\textsc{Monoide}}
\newcommand{\SET}{\textsc{Set}}
\newcommand{\MAN}{\textsc{Man}}
\newcommand{\GRUPPEN}{\textsc{Gruppen}}
\newcommand{\ABELGRUPPEN}{\textsc{Abel.Gruppen}}
\newcommand{\ABEL}{\textsc{Abel}}
\newcommand{\KAT}{\textsc{Kat}}
\newcommand{\FUN}{\textsc{Fun}}
\newcommand{\SIMP}{\textsc{Simp}}
\newcommand{\VEKT}{\textsc{Vekt}}
\newcommand{\CH}{\textsc{Ch}}
\newcommand{\CSTARUN}{C^*\text{-}\textsc{Alg}^{\raisemath{-2.5pt}{1}}}
\newcommand{\CSTAR}{C^*\text{-}\textsc{Alg}}
\newcommand{\AB}{\textsc{Ab}}
% ======================================================================================
% ======================================================================================



% -- theorem packages
% ======================================================================================
\usepackage{amsthm}
\usepackage{thmtools,thm-restate}
\usepackage{mdframed}
\renewcommand{\listtheoremname}{Übersicht aller Aussagen}
\usepackage{bookmark}
\bookmarksetup{open,numbered}
\makeatletter
\newcommand*{\theorembookmark}{%
  \bookmark[
    dest=\@currentHref,
    rellevel=1,
    keeplevel,
  ]{%
    \thmt@thmname\space\csname the\thmt@envname\endcsname
    \ifx\thmt@shortoptarg\@empty
    \else
      \space(\thmt@shortoptarg)%
    \fi
  }%
}   
\makeatother
% ======================================================================================

% -- Definition der einzelnen Theorem-Umgebungen
% ======================================================================================
\declaretheoremstyle[%
	headfont=\sffamily\bfseries,
	notefont=\normalfont\sffamily\scshape,
	bodyfont=\normalfont,
	headformat=\NUMBER\ \NAME\NOTE,
	headpunct=.,
	postheadspace=1em,
	spaceabove=15pt,spacebelow=10pt,
	shaded={bgcolor=gray!20},
	postheadhook=\theorembookmark]%
{mainstyle}
\declaretheoremstyle[%
	headfont=\sffamily\bfseries,
	notefont=\normalfont\sffamily\scshape,
	bodyfont=\normalfont,
	headformat=\NUMBER\ \NAME\NOTE,
	headpunct=.,
	postheadspace=1em,
	spaceabove=15pt,spacebelow=10pt,
	shaded={bgcolor=fb10_blue!20},
	postheadhook=\theorembookmark]%
{mainstyle_blue}
\declaretheoremstyle[%
	headfont=\sffamily\bfseries,
	notefont=\normalfont\sffamily\scshape,
	bodyfont=\normalfont,
	headformat=\NUMBER\ \NAME\NOTE,
	headpunct=.,
	postheadspace=1em,
	spaceabove=15pt,spacebelow=10pt,
	postheadhook=\theorembookmark]%
{mainstyle_unshaded}
\declaretheoremstyle[%
	headfont=\sffamily\bfseries,
	notefont=\normalfont\sffamily\scshape,
	bodyfont=\normalfont,
	headformat=\NUMBER\NAME\NOTE,
	headpunct=.,
	postheadspace=1em,
	spaceabove=15pt,spacebelow=10pt,
	% shaded={bgcolor=gray!20},
	postheadhook=\theorembookmark]%
{mainstyle_unnumbered}
\declaretheoremstyle[%
	headfont=\sffamily\bfseries,
	notefont=\normalfont\sffamily\scshape,
	bodyfont=\normalfont,
	headformat=swapnumber,
	headpunct=.,
	postheadspace=1em,
	spaceabove=15pt,spacebelow=10pt,
	shaded={bgcolor=gray!20},
	postheadhook=\theorembookmark,
	qed=\qedsymbol]%
{mainstyleB}
\declaretheoremstyle[%
	headfont=\bfseries\scshape,
	bodyfont=\normalfont,
	headpunct=:,
	postheadspace=1em,
	spacebelow=12pt,spaceabove=2pt,
	qed=\qedsymbol]%
{beweise}
\declaretheoremstyle[%
	headfont=\bfseries\scshape,
	bodyfont=\normalfont,
	headpunct=:,
	postheadspace=1em,
	spacebelow=12pt,spaceabove=2pt]%
{beweisskizze}
\declaretheoremstyle[%
	headfont=\sffamily\bfseries,
	bodyfont=\normalfont,
	headpunct=:,
	postheadspace=1em,
	spacebelow=10pt,spaceabove=10pt]%
{bemerkungen}
\declaretheorem[name=Definition,parent=section,style=mainstyle_blue]{definition}
\declaretheorem[name=Definition \& Proposition,refname=Proposition,sharenumber=definition,style=mainstyle_blue]{definitionP}
\declaretheorem[name=Definition,numbered=no,style=mainstyle_unnumbered]{definition*}
\declaretheorem[name=Theorem,sharenumber=definition,style=mainstyle]{theorem}
\declaretheorem[name=Theorem,numbered=no,style=mainstyle_unnumbered]{theorem*}
\declaretheorem[name=Proposition,sharenumber=definition,style=mainstyle,refname=Proposition]{proposition}
\declaretheorem[name=Lemma,sharenumber=definition,style=mainstyle]{lemma}
\declaretheorem[name=Satz,sharenumber=definition,style=mainstyle,refname=Satz]{satz}
\declaretheorem[name=Satz,sharenumber=definition,style=mainstyle_unshaded]{satzUnshaded}
\declaretheorem[name=Definition,sharenumber=definition,style=mainstyle_unshaded]{definitionUnshaded}
\declaretheorem[name=Satz,numbered=no,style=mainstyle_unnumbered]{satz*}
\declaretheorem[name=Korollar,sharenumber=definition,style=mainstyle,refname=Korollar]{korollar}
\declaretheorem[name=Korollar,sharenumber=definition,style=mainstyleB,refname=Korollar]{korollarB}
\declaretheorem[name=Frage,numbered=no,style=mainstyle_unnumbered]{frage}
\declaretheorem[name=Frage,sharenumber=definition,style=mainstyle_unshaded]{frageA}
\declaretheorem[name=Erinnerung,sharenumber=definition,style=mainstyle_unshaded]{erinnerungA}
\declaretheorem[name=Ausblick,sharenumber=definition,style=mainstyle_unshaded]{ausblick}
\declaretheorem[name=Konvention,sharenumber=definition,style=mainstyle]{konvention}
\declaretheorem[name=Notation,sharenumber=definition,style=mainstyle_unshaded]{notation}
\declaretheorem[name=Bemerkung,sharenumber=definition,style=mainstyle_unshaded,refname=Bemerkung]{bemerkung}
\declaretheorem[name=Bemerkung,numbered=no,style=mainstyle_unnumbered]{bemerkung*}
\declaretheorem[name=Beispiel,sharenumber=definition,style=mainstyle_unshaded,refname=Beispiel]{beispiel}
\declaretheorem[name=Beispiel,numbered=no,style=mainstyle_unnumbered]{beispiel*}
\declaretheorem[name=Exkurs,numbered=no,style=mainstyle_unnumbered]{exkurs*}
\declaretheorem[name=Beweis,numbered=no,style=beweise]{beweis}
\declaretheorem[name=Übung,numbered=no,style=bemerkungen]{uebung}
\declaretheorem[name=Erinnerung,numbered=no,style=bemerkungen]{erinnerung}

% english versions
\declaretheorem[name=Remark,sharenumber=definition,style=mainstyle_unshaded]{remark}
\declaretheorem[name=Remark,numbered=no,style=mainstyle_unnumbered]{remark*}
\declaretheorem[name=Example,sharenumber=definition,style=mainstyle_unshaded]{example}
\declaretheorem[name=Corollary,sharenumber=definition,style=mainstyle]{corollary}
\let\proof\relax
\declaretheorem[name=Proof,numbered=no,style=beweise]{proof}
\declaretheorem[name=Sketch of Proof,numbered=no,style=beweisskizze]{sketch}
% ======================================================================================

%--Inhaltsverzeichnis
% ======================================================================================
\usepackage[tocindentauto]{tocstyle}
\usetocstyle{KOMAlike}
% ======================================================================================

%-- Dinge, die erst am Ende getan werden dürfen
% ======================================================================================
\shorthandon{"}
\usepackage{ellipsis}
% ======================================================================================


\newcommand{\fach}{Lie-Gruppen}
\newcommand{\semester}{Sose 2016}
\newcommand{\homepage}{http://wwwmath.uni-muenster.de/42/arbeitsgruppen/ag-differentialgeometrie/prof-dr-christoph-boehm/vorlesung-liegruppen/}

\newcommand{\prof}{Prof.\ Dr.\ Christoph Böhm}
\publishers{\scalebox{10}{\Huge$\mathrm{Lie}$}}
\providecommand{\verfasser}{Jannes Bantje}

%--Konfiguration von scrheadings
\setheadsepline{1pt}[\color{light_gray}]
\pagestyle{scrheadings}
%\fancyhf[H, F]{}
\clearscrheadfoot

\providecommand{\shortFach}{\fach}
\lehead{\includegraphics[height=0.6 cm,keepaspectratio]{../!config/Bilder/Logo_WWU_Muenster_light_gray.pdf}}
\rehead{\rule{0cm}{0.6cm}\footnotesize\sffamily\color{light_gray}\verfasser{} -- Mitschrift \shortFach}
\lohead{\rule{0cm}{0.6cm}\footnotesize\sffamily\color{light_gray}Stand: \today \; \thistime[:]}
\rohead{\includegraphics[height=0.6 cm,keepaspectratio]{../!config/Bilder/fb10logo_gray.pdf}}


\ofoot[{ \color{dark_gray} \LARGE \sffamily \thepage}]{{ \color{dark_gray} \LARGE \sffamily \thepage}} %hier wir auch der plain Stil bearbeitet!
\automark{section}
\ifoot{ \color{dark_gray} \small \leftmark}

%--Metadaten
\providecommand{\mail}{j.bantje@wwu.de}
\author{\verfasser}
\titlehead{\includegraphics[height=1.5cm, keepaspectratio]{../!config/Bilder/Logo_WWU_Muenster.pdf}%
\hfill \includegraphics[height=1.3cm, keepaspectratio]{../!config/Bilder/fb10logo.pdf}}
\title{Skript \fach}
\subtitle{Mitschrift der Vorlesung  \enquote{\fach} von \prof}

\KOMAoptions{DIV=last}


\begin{document}
\pagenumbering{Roman}
\maketitle
\begin{abstract}
\section*{Aktuelle Version verfügbar bei}
\newcommand{\dieBreite}{11cm}
\begin{minipage}{4cm}
	\qrcode[height=3.3cm, version=6]{https://gitlab.com/JaMeZ-B/LaTeX-WWU}
\end{minipage}
\hfill
\begin{minipage}{\dieBreite}
	% \includegraphics[height=0.6cm, keepaspectratio]{../!config/Bilder/wm_no_bg.pdf}
	\includegraphics[height=0.8cm, keepaspectratio]{../!config/Bilder/wm_no_bg.pdf}\\
	\url{https://gitlab.com/JaMeZ-B/LaTeX-WWU} \smallskip\\
	Das zentrale Repository des \enquote{\LaTeX-WWU}-Projekts befindet sich auf der Plattform GitLab.com.
	Neben der Koordination aller Beteiligten werden über diesen Dienst auch die PDFs gebaut, die in der Readme verlinkt sind.
\end{minipage}\\[1cm]
\begin{minipage}{4cm}
	\qrcode[height=3.3cm, version=6]{https://github.com/JaMeZ-B/latex-wwu}
\end{minipage}
\hfill
\begin{minipage}{\dieBreite}
	\includegraphics[height=0.6cm, keepaspectratio]{../!config/Bilder/github_octo.pdf}
	\includegraphics[height=0.6cm, keepaspectratio]{../!config/Bilder/GitHub_Logo.pdf}\\
	\url{https://github.com/JaMeZ-B/latex-wwu} \smallskip\\
	Die Entwicklung des \enquote{\LaTeX-WWU}-Projekts hat ursprünglich auf GitHub stattgefunden, ist mittlerweile aber zu GitLab gewechselt.
	Das GitHub-Repository wird stündlich automatisch aktualisiert, Merge-Requests werden aber nicht mehr entgegengenommen.
\end{minipage}\\[1cm]
% \begin{minipage}{4cm}
% 	\qrcode[height=3.3cm, version=6]{https://uni-muenster.sciebo.de/public.php?service=files&t=965ae79080a473eb5b6d927d7d8b0462}
% \end{minipage}
% \hfill
% \begin{minipage}{\dieBreite}
% 	\raisebox{-2pt}{\includegraphics[height=0.6cm, keepaspectratio]{../!config/Bilder/sciebo_logo.pdf}}
% 	\resizebox{!}{0.5cm}{\large \sffamily\textbf{sciebo}} {\sffamily\large die Campuscloud} \\
% 	\resizebox{\dieBreite}{!}{\footnotesize\url{https://uni-muenster.sciebo.de/public.php?service=files&t=965ae79080a473eb5b6d927d7d8b0462}}\smallskip\\
% 	Sciebo ist ein Dropbox-Ersatz der Hochschulen in NRW, der von der Uni Münster in leitender Position auf Basis der OpenSource-Software Owncloud aufgebaut wurde.
% \end{minipage}\\[1cm]
\hrule \mbox{ }\\[0.7cm]
\begin{minipage}{4cm}
	\qrcode[height=3.3cm, version=6]{\homepage}
\end{minipage}
\hfill
\begin{minipage}{\dieBreite}
	\resizebox{!}{0.5cm}{\large\sffamily\textbf{Vorlesungshomepage}}\\
	\resizebox{\dieBreite}{!}{\footnotesize\url{\homepage}}\smallskip\\
	Hier ist ein Link zur offiziellen Vorlesungshomepage.
\end{minipage}
\newpage
\section*{Vorwort --- Mitarbeit am Skript}
Dieses Dokument ist eine Mitschrift aus der Vorlesung \enquote{\fach, \semester}, gelesen von \prof. 
Der Inhalt entspricht weitestgehend dem Tafelanschrieb. 
Für die Korrektheit des Inhalts übernehme ich keinerlei Garantie! 
Für Bemerkungen und Korrekturen -- und seien es nur Rechtschreibfehler -- bin ich sehr dankbar. 
Korrekturen lassen sich prinzipiell auf drei Wegen einreichen: 
\begin{itemize}
	\item Persönliches Ansprechen in der Uni, Mails an \hrefsymmail{mailto:\mail}{\mail} (gerne auch mit annotieren PDFs) oder Kommentare auf \url{https://gitlab.com/JaMeZ-B/LaTeX-WWU}.
	\item \emph{Direktes} Mitarbeiten am Skript: Den Quellcode poste ich auf GitLab (siehe oben), also stehen vielfältige Möglichkeiten der Zusammenarbeit zur Verfügung:
	Zum Beispiel durch Kommentare am Code über die Website und die Kombination Fork und Merge-Request. 
	Wer sich verdient macht oder ein Skript zu einer Vorlesung, die ich nicht besuche, beisteuern will, dem gewähre ich gerne auch Schreibzugriff.
	
	Beachten sollte man dabei, dass dazu ein Account bei \url{gitlab.com} notwendig ist, der allerdings ohne Angabe von persönlichen Daten angelegt werden kann. 
	Wer bei GitLab (bzw. dem zugrunde liegenden Open-Source-Programm \enquote{\texttt{git}}) -- verständlicherweise -- Hilfe beim Einstieg braucht, dem helfe ich gerne weiter. 
	Es gibt aber auch zahlreiche empfehlenswerte Tutorials im Internet.\footnote{zB. \url{https://try.github.io/levels/1/challenges/1}, ist auf Englisch, aber dafür interaktiv}
	\item \emph{Indirektes} Mitarbeiten: \TeX-Dateien per Mail verschicken. 
	
	Dies ist nur dann sinnvoll, wenn man einen ganzen Abschnitt ändern möchte (zB. einen alternativen Beweis geben), da ich die Änderungen dann per Hand einbauen muss! Ich freue mich aber auch über solche Beiträge!
\end{itemize}
\section*{Literatur}
\begin{itemize}
	\item \citetitle{Ziller} von Wolfgang \citeauthor{Ziller} \cite{Ziller}
	\item \citetitle{Berndt} von Jürgen \citeauthor{Berndt} \cite{Berndt} 
\end{itemize}
\end{abstract}

\tableofcontents
\cleardoubleoddemptypage

\pagenumbering{arabic}
\setcounter{page}{1}
\setcounter{footnote}{0}

\chapter{Grundlagen} % (fold)
\label{cha:grundlagen}

% chapter grundlagen (end)

\section{Liegruppen und Liealgebren} % (fold)
\label{sec:1}
\begin{definition}[label=def:111,{name=[Liegruppe]}]
	Eine \Index{Liegruppe} st eine abstrakte Gruppe, welche zusätzlich eine $n$-dimensionale, differenzierbare Mannigfaltigkeit ist, sodass die beiden Abbildungen 
	$m \colon G \times G \to G$, $(g,h) \mapsto g \cdot h$ und $^{-1} \colon G \to G$, $g \mapsto g^{-1}$ differenzierbar sind.
\end{definition}

Beispiele für Liegruppen: 
\begin{itemize}
	\item Die $1$-Sphäre $S^1 = \set*{e^{i \varphi} \in \mathbb{C} \given \varphi \in \mathbb{R}} = \set[\big]{z \in \mathbb{C} \given \abs*{z}=1} \simeq \set*{\begin{psmallmatrix*}[r]
		\cos \varphi & - \sin \varphi \\
		\sin \varphi & \cos \varphi
	\end{psmallmatrix*} \given \varphi \in \mathbb{R}} = \SO(\mathbb{R})$
	\item Der $n$-Torus $T^n = S^1 \times \ldots \times S^1 $ ($n$-mal)
	\item Sei $r \in \mathbb{R}$ und $S^1_r := \set*{\enbrace{e^{i \varphi}, e^{i r \varphi}} \in S^1 \times S^1 \given \varphi \in \mathbb{R}}$.
	Für $r \in \mathbb{Q}$ ist $S^1_r$ eine $1$-dimensionale Untermannigfaltigkeit des Torus $T^2= S^1 \times S^1$.
	Für $r \in \mathbb{R} \setminus \mathbb{Q}$ ist $S^1_r$ \emph{keine} eingebettete Untermannigfaltigkeit, denn es gilt 
	\[
		\overline{S^1_r} = T^2
	\]
	\item Die $3$-Sphäre $S^3 = \set[\big]{q \in \mathbb{H} \given \abs*{q}=1}$, wobei $\mathbb{H}$ der Schiefkörper der \Index{Quaternionen}\footnote{siehe \url{https://de.wikipedia.org/wiki/Quaternion}} ist.
	
	Wir setzen $\mathbb{R}^3 := \im \mathbb{H} = \Span_\mathbb{R} (i,j,k)$.
	Dann ist $\mathbb{R}^3$ invariant unter der adjungierten Darstellung von $S^3$, das heißt für alle $x \in \im \mathbb{H}$ und alle $q \in S^3$ gilt:
	\[
		\Ad(q)(x) := q \cdot x \cdot \Underbrace{q^{-1}}{=\overline{q}} \stackrel{!}{\in} \mathbb{R}^3 \qquad \text{ für } x \in \mathbb{R}^3 
	\]
	Man kann ferner zeigen, dass $\Ad(q)|_{\mathbb{R}^3} \subset SO(3)$.
	Es gilt sogar $\Ad(S^3) = SO(3)$ und $\pi \colon S^3 \to SO(3)$, $q \mapsto \Ad(q)|_{\mathbb{R}^3}$ ist eine $2$-fache Überlagerung.
	Dies folgt mit $\Ad(q)=\Ad(-q)$.
	\item \Index{Orthogonale Gruppe} $\On(n) = \set*{A \in \Mat(n,\mathbb{R}) \given A A^T = E_n}$, $\SO(n) = \set[\big]{A \in \On(n) \given \det A=1}$
	\item \Index{Unitäre Gruppe} $\Un(n) = \set[\big]{A \in \Mat (n,\mathbb{C}) \given A A^* =E_n}$, $\SU(n) = \set[\big]{A \in \Un(n) \given \det A=1}$
	\item \Index{Allgemeine lineare Gruppe}: $\GL(n,\mathbb{K}) = \set[\big]{A \in \Mat(n,\mathbb{K}) \given A \text{ invertierbar}}$. \Index{Spezielle lineare Gruppe}:
	$\SL(n,\mathbb{K}) = \set[\big]{A \in \GL(n,\mathbb{K}) \given \det A =1}$
	\item Sind $G_1$ und $G_2$ Liegruppen, so ist auch $G_1 \times G_2$ eine Liegruppe.
	\item Ist $G$ eine Liegruppe, so ist 
	\[
		G_0 :=  \set[\big]{g \in G \given \exists c \colon [0,1] \to G \text{ stetig mit } c(0)=e, c(1)=g}
	\]
	ein Liegruppe, die Zusammenhangskomponente der Eins.
	Beispiel: $\On(n)=\SO(n) \mathbin{\dot{\cup}} I \cdot \SO(n)$, wobei $I$ eine Spiegelung ist.
	Tatsächlich hat $\On(n)$ genau zwei Zusammenhangskomponenten, für $n=2$:
	\[
		\On(2) = \SO(2) \mathbin{\dot{\cup}} \set*{\begin{pmatrix*}[r]
			\cos \varphi & \sin \varphi \\
			\sin \varphi & - \cos \varphi
		\end{pmatrix*} \given \varphi \in \mathbb{R}}
	\]
\end{itemize}

\begin{definition}[{name=[Lieklammer]}]
	Eine \Index{Liealgebra} über $\mathbb{R}$ oder $\mathbb{C}$ ist ein $\mathbb{K}$-Vektorraum $V$ zusammen mit einer bilinearen, schiefsymmetrischen Abbildung 
	\(
		[\cdot ,\cdot ] \colon V \times V \to V
	\)
	welche die Jacobi-Identität 
	\[
		\benbrace[\big]{X, \benbrace*{Y,Z}} + \benbrace[\big]{Z, \benbrace*{X,Y}} + \benbrace[\big]{Y, \benbrace*{Z,X}} =0
	\]
	für alle $X,Y,Z \in V$ erfüllt.
	Man nennt $\benbrace*{\cdot ,\cdot }$ eine \Index{Lieklammer}. 
\end{definition}

Wir werden zeigen, dass für eine Liegruppe $G$ der Tangentialraum $\Tmap_e G$ auf natürliche Art und Weise eine Liealgebra ist.
Dazu ein wenig Notation: Für $g \in G$ betrachten wir Links- und Rechtsmultiplikation
\begin{align}
	\SwapAboveDisplaySkip
	L_g \colon G \longrightarrow G \quad , \qquad &h \longmapsto g \cdot h \\
	R_g \colon G \longrightarrow G \quad , \qquad &h \longmapsto h \cdot g
\end{align}
Die Links-, sowie die Rechtsmultiplikation sind Diffeomorphismen. 

\begin{definition}[{name=[linksinvariantes Vektorfeld]}]
	Ein Vektorfeld $X$ auf $G$ nennt man \Index{linksinvariant}, falls für alle $g,h \in G$ gilt
	\[
		X(g \cdot h) = \mathd(L_g)_h \cdot X(h)
	\]
\end{definition}

Offenbar ist ein linksinvariantes Vektorfeld $X$ schon eindeutig durch seinen Wert $X(e)$ bestimmt, denn 
\[
	X(g) = X(g \cdot e) = \mathd(L_g)_e \cdot X(e)
\]
Wir schreiben daher auch $X=X_v$ mit $X(e)=v \in \Tmap_e G$.

\begin{lemma}[label=lem:114,{name=[Linksinvariante Vektorfelder sind differenzierbar]}]
	Linksinvariante Vektorfelder sind differenzierbar.
\end{lemma}
\begin{beweis}
	Nach \autoref{def:111} ist die Multiplikation $m \colon G \times G \to G$ differenzierbar, somit auch $\mathd m \colon \Tmap G \times \Tmap G \to\Tmap G$ und damit für alle $v \in \Tmap_e G$ auch $\tilde{X}_v \colon G \to \Tmap G$, $g \mapsto (\mathd m)_{(g,e)}(0,v)$.
	Wir zeigen nun $\tilde{X}_v=X_v$.
	Sei $\enbrace[\big]{g(t)}_{t \in (-1,1)}$ eine differenzierbare Kurve in $G$ mit $g(0)=e$ und $g'(0)=v$.
	Dann gilt 
	\[
		(\mathd m)_{(g,e)}(0,v) = \diffd{}{t}\Big|_{t=0} m \enbrace[\big]{g,g(t)} = \diffd{}{t}\Big|_{t=0} L_g\enbrace[\big]{g(t)} = \mathd(L_g)_{g(0)} g'(0)
		 =\mathd(L_g)_e v = X_v(g) \qedhere
	\]
\end{beweis}

Sei $\mathfrak{g}$ die Menge der linksinvarianten Vektorfelder auf $G$.

\begin{korollar}
	Die Abbildung $L \colon \mathfrak{g} \to \Tmap_e G, X \mapsto X(e)$ ist ein linearer Isomorphismus.
\end{korollar}
\begin{beweis}
	Die Injektivität ist klar.
	Für die Surjektivität sei $v \in \Tmap_e G$ ein Tangentialvektor.
	Nach \autoref{lem:114} ist $X_v$ ein glattes Vektorfeld auf $G$, es ist weiter linksinvariant, da
	\begin{align}
		X_v(g \cdot h) = \mathd(L_{g \cdot h})_e v = \mathd(L_g)_h \cdot (\mathd L_h)_e \cdot v = \mathd (L_g)_h \cdot X_v(h)
	\end{align}
	Ein Vektorfeld ist linksinvariant genau dann, wenn $X$ für alle $g \in G$ $L_g$-verwandt zu sich selbst ist, das heißt
	\[
		X \enbrace[\big]{L_g(h)} = (\mathd L_g)_h X(h) \qquad  \forall h \in G
	\]
	Nach Diffgeo I ist die Lieklammer von $L_g$-verwandten Vektorfeldern wieder $L_g$-verwandt.
	Die Lieklammer linksinvarianter Vektorfelder ist also wieder ein linksinvariantes Vektorfeld
\end{beweis}

\begin{definition}[{name=[Lieklammer der Liealgebra]}]
	Wir bezeichnen mit 
	\[
		\benbrace*{\cdot ,\cdot } \colon \Tmap_e G \times \Tmap_e G \to \Tmap_e G \quad ,\quad  (v,w) \longmapsto \benbrace*{X_v,X_w}(e)
	\] 
	die \Index{Lieklammer} von $\Tmap_e G = \mathfrak{g}$.
	Da die Lieklammer differenzierbarer Vektorfelder die Jacobi-Identität erfüllt, ist $\benbrace*{\cdot ,\cdot }$ eine Lieklammer auf $\Tmap_e G$.
\end{definition}

Sei $G=\GL(n,\mathbb{K})$, $\mathbb{K}=\mathbb{R}$ ($\mathbb{K}=\mathbb{C}$).
Da die Determinante $\det \colon G \to \mathbb{K}$ stetig ist, ist 
\[
	{\det}^{-1} \enbrace[\big]{\mathbb{R}\setminus \set*{0}} = \GL(n,\mathbb{K}) 
\]
eine offene Teilmenge des euklidischen Vektorraums $\Mat(n,\mathbb{K})$ und somit eine differenzierbare Mannigfaltigkeit.
Wir bezeichnen mit $x_{ij} \colon G \to \mathbb{K}$ die Abbildung $(a_{kl})_{k,l} \mapsto a_{ij}$ für festes $(i,j)$ mit $1\le i,j\le n$.
Da $x_{ij}$ linear ist, gilt für $A \in G$
\[
	(X_v)(x_{ij})(A) = (\mathd x_{ij})_A X_v(A) = x_{ij}(A \cdot v) = (A v)_{ij}
\]
Somit ist
\begin{align}
	\benbrace[\big]{X_v,X_w}(x_{ij})(e) = X_v(e) \enbrace[\big]{X_w(x_{ij})} - X_w(e) \enbrace[\big]{X_v(x_{ij})} &= v \enbrace[\Big]{\Underbrace{X_w(x_{ij})}{\mathclap{A \mapsto (A \cdot w)_{ij}}}} - w \enbrace*{X_v(x_{ij})} \\
	&= (v w- w v)_{ij}
\end{align}
Für alle differenzierbaren Vektorfelder $X$ und Koordinaten $(x_1, \ldots ,x_n)=x$  gilt in einem Kartengebiet $U$
\[
	X = \sum_{i=1}^{n} X(x_i) \cdot \diff{}{x_i}
\]
Somit gilt $\benbrace*{X_v,X_w} =X(vw -wv)$ und daher $\benbrace*{v,w}_{\mathfrak{gl}(n,\mathbb{R})}=vw -wv$
% subsection 11 (end)

\section{Lieuntergruppen und Homomorphismen von Liegruppen} % (fold)
\label{sec:12}

\begin{definition}[{name=[Liegruppen- und Liealgebrenhomomorphismen]}]
	Seien $H,G$ Liegruppen.
	\begin{enumerate}[1)]
		\item Eine Abbildung $\Phi \colon H \to G$ nennt man \Index{Liegruppenhomomorphismus}, falls $\Phi$ differenzierbar ist und falls $\Phi$ ein Gruppenhomomorphismus ist.
		\item Einen Liegruppenhomomorphismus $\Phi \colon H \to G$ nennt man \Index{Liegruppenisomorphismus}, falls $\Phi$ bijektiv und $\Phi^{-1}$ differenzierbar ist. ($\Phi^{-1}$ ist auch Gruppenhomomorphismus)
		\item \bet{Liealgebrenhomomorphismen}\index{Liealgebrenhomomorphismus} und Liealgebrenisomorphismen werden entsprechend definiert:  $L \colon \mathfrak{h} \to \mathfrak{g}$ linear mit
		\[
			L \cdot \benbrace*{v,w}_{\mathfrak{h}} = \benbrace[\big]{L \cdot v, L \cdot w}_\mathfrak{g}
		\]
	\end{enumerate}
\end{definition}

Ist $\Phi \colon H \to G$ ein Liegruppenhomomorphismus, so kommutiert das folgende Diagramm
\[
	\begin{tikzcd}
		H \rar["\Phi"] \dar["L_h"] & G \dar["L_{\Phi(h)}"] \\
		H \rar["\Phi"] & G
	\end{tikzcd}
\]
da
\begin{align}
	(\Phi \circ L_h)(\tilde{h}) = \Phi \enbrace*{L_h(\tilde{h})} = \Phi(h \cdot \tilde{h}) = \Phi(h) \cdot \Phi(\tilde{h}) = \enbrace*{L_{\Phi(h)} \circ \Phi}(\tilde{h})
\end{align}
Einen Homomorphismus $\Phi \colon G \to \GL(n,\mathbb{K})$ nennt man \bet{reelle}\index{reelle Darstellung} bzw. \Index{komplexe Darstellung} von $G$. ($\Phi_0(G)\equiv E_n$)
Ist $\Phi \colon G \to \GL(n,\mathbb{K})$ injektiv, so besitzt die Gruppe $G$ die \Index{treue Darstellung} $\Phi(G)$ in $\GL(n,\mathbb{K})$.

\begin{lemma}[label=lem:122]
	Sei $\Phi \colon H \to G$ ein Liegruppenhomomorphismus.
	Dann ist $(\mathd\Phi)_e \colon \Tmap_e H \to \Tmap_e G$ ein Liealgebrenhomomorphismus.
\end{lemma}
\begin{beweis}
	Seien $v_1,v_2 \in \Tmap_e H$ und $w_i := (\mathd \Phi)_e \cdot v_i$ für $i=1,2$.\marginnote{$\mathd \Phi=: L$}
	Zu zeigen
	\[
		(\mathd \Phi)_e \benbrace*{X_{v_1}, X_{v_2}} (e) = \benbrace*{X_{w_1},X_{w_2}}(e)
	\]
	Wir zeigen, dass $X_{v_1}$ und $X_{w_1}$ $\Phi$-verwandt sind:
	\begin{align}
		(\mathd \Phi)_h X_{v_1}(h) = \enbrace*{\mathd \Phi}_h \cdot (\mathd L)_e \cdot v_1 &= \mathd \enbrace*{\Phi \circ L_h}_e \cdot v_1 \\
		&= \mathd \enbrace*{L_{\Phi(h)} \circ \Phi}_e \cdot v_1 \\
		&= \mathd \enbrace*{L_{L_{\Phi(h)}}} \cdot \enbrace*{\mathd \Phi}_e \cdot v_1 \\
		&= \mathd \enbrace*{L_{\Phi(h)}} \cdot w_1 \\
		&= X_{w_1}(\Phi(h))
	\end{align}
	Somit sind $X_{v_1}$ und $X_{w_1}$ $\Phi$-verwandt , analog für $v_2$ und $w_2$.
	Mit DiffGeo \RM{1} sind dann auch $\benbrace*{X_{v_1}, X_{v_2}}$ und $\benbrace*{X_{w_1},X_{w_2}}$ $\Phi$-verwandt, das heißt
	\[
		(\mathd \Phi)_h \benbrace*{X_{v_1},X_{v_2}}(h) = \benbrace*{X_{w_1}, X_{w_2}} \enbrace*{\Phi(h)}
	\]
	Mit $h=e$ folgt die Behauptung.
\end{beweis}

Sei $G$ eine Liegruppe und $i_g \colon G \to G$ gegeben durch $h \mapsto g \cdot h \cdot g^{-1}$ die \Index{Konjugation} mit $g$.
Dann ist $i_g$ ein Liegruppenisomorphismus
\begin{itemize}
	\item $i_g(h) = \enbrace*{l_g \circ R_{g^{-1}}}(h) \Rightarrow$ differenzierbar
	\item Die Gruppenhomomorphismuseigenschaft ist klar. 
	\item $(i_g)^{-1}= i_{g^{-1}}$
\end{itemize}
Somit ist $\Ad(g) := \enbrace*{\mathd i_g}_e \colon \mathfrak{g} \to \mathfrak{g}$ ein Liealgebrenisomorphismus.
Ist $G = \GL(n,\mathbb{K})$, so kann man $\Ad(g)$ einfach berechnen:
Sei $(h(t))_{t \in (-\varepsilon,\varepsilon)}$ eine differenzierbare Kurve in $G$ mit $h(0)= E_n$, $h'(0)=v$.
Dann gilt
\begin{align}
	\Ad(g)(v)= (\mathd i_g)_e \cdot h'(0) = \diffd{}{t}\Big|_{t=0} i_g (h(t)) = \diffd{}{t}\Big|_{t=0} g \cdot h(t) \cdot g^{-1}
	&= g \cdot \diffd{}{t}\Big|_{t=0} h(t) \cdot g^{-1}\\
	&= g \cdot v \cdot g^{-1}
\end{align}
Ferner ist $\Ad \colon G \to \GL\enbrace*{\Tmap_e G}$, $g \mapsto \Ad(g)$ ein Liegruppenhomomorphismus:
\begin{align}
	\Ad(g_1 \cdot g_2) = (\mathd i_{g_1 \cdot g_2})(e) = (\mathd i_{g_1})_e \cdot (\mathd i_{g_2})_e = \Ad(g_1) \circ \Ad(g_2)
\end{align}
$\Ad \colon G \to \GL(\mathfrak{g})$ heißt die \Index{adjungierte Darstellung} von $G$.

\begin{definition}[{name=[{Lieuntergruppe}]}]
	Sei $G$ eine Liegruppe.
	\begin{enumerate}[1)]
		\item Mann nennt eine Teilmenge $H \subseteq G$ eine \Index{Lieuntergruppe}, falls 
		\begin{itemize}
			\item $H$ eine Untergruppe von $G$ ist
			\item $H$ ist Liegruppe, wenn $H$ mit der von $G$ induzierten Multiplikation und Inversenbildung versehen wird.
			\item Die Inklusionsabbildung $i \colon H \hookrightarrow G$ ist eine Immersion \marginnote{Immersion: injektives Differential}\\ 
			(nicht notwendigerweise eine Einbettung!).
		\end{itemize}
		\item Mann nennt einen Teilraum $\mathfrak{h} \subseteq \mathfrak{g}$ Lieunteralgebra, falls $\benbrace*{x,y}_\mathfrak{g} \in H$ für alle $x,y \in \mathfrak{h}$ gilt.
	\end{enumerate}
\end{definition}
 
Die Gruppe $S^1_r = \set*{\enbrace*{e^{i \varphi}, e^{i r \varphi}} \given \varphi \in \mathbb{R}}$ ist eine Lieuntergruppe von $T^2=S^1 \times S^1$.
Für $r \in \mathbb{Q}$ ist $S^1_r$ eine (eingebettete) Untermannigfaltigkeit; für $r \in \mathbb{R}\setminus \mathbb{Q}$ ist $S^1_r$ \enquote{nur} eine immersierte Untermannigfaltigkeit, jedoch nicht eingebettet.

$\SO(n) \subset \GL(n,\mathbb{R})$ ist eine Lieuntergruppe mit Liealgebra 
\[
	\mathfrak{so}(n) := \Tmap_e \SO(n) = \set*{X \in \Mat(n,\mathbb{R}) \given X^T=-X}
\]
Sei dazu $(A(t))_{t \in (-\varepsilon,\varepsilon)}$ eine differenzierbare Kurve mit $A(0)=E_n$ und $A'(0)=X$.
Aus $A(t)^T \cdot A(t) \equiv E_n$ folgt durch Differenzieren $X^T + X =0$.
Sei umgekehrt $X$ eine schiefsymmetrische Matrix.
Setze dann 
\[
	A(t) =\exp(t X) = \sum_{k=0}^{\infty} \frac{(tX)^k}{k!} 
\]
Behauptung: $A(t) \in \SO(n)$.
\begin{align}
	A(t)^T \cdot A(t) = \enbrace*{\sum_{k=0}^{\infty} \frac{(tX)^k}{k!} }^T \cdot \sum_{k=0}^{\infty} \frac{(tX)^k}{k!} &\stackrel{X^T=-X}{= }\sum_{k=0}^{\infty} \frac{(-tX)^k}{k!}  \cdot \sum_{k=0}^{\infty} \frac{(tX)^k}{k!} \\
	&= \exp(x-x) = \exp(0) =E_n
\end{align}
die folgt mit $\exp(x+y)= \exp(x) \exp(y)$, da $[x,y]=0$.

\begin{lemma}
	Sei $G$ eine Liegruppe und $H \subset G$ eine Lieuntergruppe.
	Dann ist $\mathfrak{h} = \Tmap_e H$ eine Lieunteralgebra von $\mathfrak{g} = \Tmap_e G$.
\end{lemma}
\begin{beweis}
	Sei $i \colon H \to G$ die Inklusionsabbildung.
	Dann ist $i$ ein Liegruppenhomomorphismus.
	Nach \autoref{lem:122} ist somit $(\mathd i)_e = \id_{\Tmap_e H} \colon \Tmap_e H \to \Tmap_e G$ ein Liealgebrenhomomorphismus, das heißt $\Tmap_e H = \mathfrak{h}$ ein Lieunteralgebra von $\mathfrak{g}$.
\end{beweis}

Wir erinnern daran, dass die Zusammenhangskomponente $G_0$ der \enquote{Eins} $e \in G$ einer Liegruppe wieder eine Liegruppe ist.
Beispiel: 
\[
	\On(n)= \SO(n) \mathbin{\dot{\cup}} \set*{ \begin{psmallmatrix}
		-1 & & &\\
		& 1 &  &\\
		& & \ddots & \\
		& & & 1
	\end{psmallmatrix} \cdot \SO(n)}
\]

\begin{satz}[label=satz:125]
	Sei $G$ eine Liegruppe und $\mathfrak{h}$ eine Unteralgebra von $\mathfrak{g}$.
	Dann existiert genau eine zusammenhängende Lieuntergruppe $H$ von $G_0$ mit Liealgebra $\mathfrak{h}$. 
\end{satz}
\begin{beweis}
	Sei $g \in G$ und $\mathcal{F}_g := (\mathd L_g) \cdot \mathfrak{h} \subset \Tmap_g G$.
	Man nennt $\mathcal{F} = \set*{\mathcal{F}_g}_{g \in G}$ eine \Index{Distribution} von Unterräumen (gleicher Dimension!) auf $G$.
	Seien $v_1, \ldots ,v_r \in \Tmap_e G$ eine Basis von $\mathfrak{h}$.
	Wir bezeichnen mit $X_{v_1}, \ldots ,X_{v_r}$ die entsprechenden linksinvarianten Vektorfelder auf $G$.
	Es gilt
	\[
		\mathcal{F}_g = \Span_\mathbb{R} \enbrace*{X_{v_1}(g), \ldots , X_{v_r}(g)}
	\]
	Die Distribution $\mathcal{F}$ ist involutiv, das heißt $\benbrace*{\mathcal{F},\mathcal{F}} \subset \mathcal{F}$, denn nach Definition ist
	\[
		\benbrace*{X_{v_i}, X_{v_i}} = X_{\benbrace*{v_i,v_i}_{\mathfrak{g}}} = X_w
	\]
	mit $w= \benbrace*{v_i,v_i} \in \mathfrak{h}$.
	Der \emph{Satz von Frobenius} (Lee: Introduction to smooth manifolds)\todo{Ref!} besagt nun, dass durch jeden Punkt $g \in G$ eine zusammenhängende, maximale Integralmannigfaltigkeit $M^r_g$ geht, das heißt $\forall \tilde{g} \in M^r_g$ gilt $\Tmap_{\tilde{g}} M^r_g = \mathcal{F}_{\tilde{g}}$.
	($M^r_g$ ist immersierte Untermannigfaltigkeit)
	Da die Distribution $\mathcal{F} = \set*{\mathcal{F}_g}_{g \in G}$ $G$-invariant ist, bildet der Diffeomorphismus $L_g \colon G \to G$ maximale Blätter (= Integralmannigfaltigkeiten) auf maximale Blätter ab.
	Wir setzen $H := M^r_e$ (wir wissen schon: $H$ ist immersierte Untermannigfaltigkeit von $G$).
	Zeige nun, dass $H$ eine Untergruppe ist:
	Da $H$ und $L_{h^{-1}}(H)$ beide $e$ enthalten und maximale Blätter sind, gilt $H=L_{h^{-1}}(H)$ für $h \in H$.
	Damit ist die Abgeschlossenheit unter Inversion gezeigt.
	Ähnlich zeigt man die Abgeschlossenheit unter Multiplikation.
	Da $m \colon G  \times G \to G$, $^{-1} \colon G \to G$ differenzierbare Abbildungen sind und die Einschränkung von differenzierbaren Abbildungen auf immersierte Untermannigfaltigkeiten wieder differenzierbare Abbildungen liefert, ist $H$ eine Lieuntergruppe von $G$.
	
	Es bleibt die Eindeutigkeit zu zeigen: 
	Sei also $\tilde{H} \subset L$ eine zusammenhängende Lieuntergruppe von $G$ mit $\Tmap_e \tilde{H} = \mathfrak{h} = \Tmap_e H$.
	Es gilt $\Tmap_{\tilde{h}} \tilde{H} = (\mathd L)_{\tilde{h}} \cdot \Tmap_e \tilde{H}= \mathcal{F}_{\tilde{h}}$ für alle $\tilde{h} \in \tilde{H}$.
	Somit ist $\tilde{H}$ Integralmannigfaltigkeit von $\mathcal{F}$ mit $e \in \tilde{H}$.
	Da $H$ maximale Integralmannigfaltigkeit ist, folgt $\tilde{H} \subseteq H$.
	Es bleibt die Gleichheit zu zeigen:
	$\tilde{H}$ ist offen in $H$.
	Annahme: Es gibt ein $h \in H \setminus \tilde{H}$ und eine Folge $(\tilde{h}_i)_{i \in \mathbb{N}}$ in $\tilde{H}$ mit $\tilde{h}_i \to h$.
	$h \in L_h(U)$, wobei $U$ eine offene Umgebung von $e$ ist.
	Damit ist $h \in L_{h_i}(U)$ für $i \ge i_0$ für geeignetes $i_0 \in \mathbb{N}$.
	Dies ist ein Widerspruch.
\end{beweis}

\begin{korollar}[label=kor:126]
	Seien $H,G$ zusammenhängende Liegruppen und $\Phi, \psi \colon H \to G$ Gruppenhomomorphismen.
	Gilt $(\mathd \Phi)_e = (\mathd \psi)_e \colon \Tmap_e H \to \Tmap_e G$, so folgt $\Phi = \psi$.
\end{korollar}
\begin{beweis}
	$H \times G$ ist eine Liegruppe mit Liealgebra $\mathfrak{h} \oplus \mathfrak{g}= \Tmap_e H \oplus \Tmap_e G$ (komponentenweise).
	Dann ist eine differenzierbare Abbildung $\Phi \colon H \to G$ genau dann ein Liegruppenhomomorphismus, wenn der Graph von $\Phi$, also $\set[\big]{(h,\Phi(h)) \given h \in H}$ eine Untergruppe von $H \times G $ ist (!).
	Ist dies der Fall, so ist 
	\[
		\Graph \enbrace[\big]{(\mathd \Phi)_e} = \set[\big]{ \enbrace*{v, (\mathd \Phi)_e \cdot v} \given v \in \mathfrak{h}}
	\]
	eine Unteralgebra von $\mathfrak{h} \oplus \mathfrak{g}$.
	Wegen $\Graph \enbrace*{(\mathd \Phi)_e} = \Graph \enbrace*{(\mathd \psi)_e}$ sind diese beiden Unteralgebren gleich und somit auch die entsprechenden eindeutigen zusammenhängenden Lieuntergruppen nach \autoref{satz:125}.
	Somit ist $\Phi =\psi$.
\end{beweis}

Wir bezeichnen mit $c_v \colon I_{\max} \to G$, $t \mapsto c_v(t)$ die maximale Integralkurve von $X_v$ mit $c_v(0)=e$, $c_v'(0)=v$.
Es gilt $c_v'(t) = X_v \enbrace*{c_v(t)}$ für alle $t$.
Wir wissen, dass $I_{\max}$ offen in $\mathbb{R}$ ist.

Beispiel:
Betrachte $\GL(n,\mathbb{R})$ und $v \in \Tmap_e G = \Mat(n,\mathbb{R})$.
\[
	X_v(A) = (\mathd L_A)_e \cdot v = A \cdot v
\]
Es gilt $c_v(t)= \exp(t \cdot v) = \sum_{k=0}^{\infty} \frac{(tv)^k}{k!}$, denn
\begin{align}
	\diffd{}{t}\Big|_{t=t_v} c_v(t) &= \diffd{}{t}\big|_{t=t_0} \Underbrace{\exp(t v)}{=\exp(t_0 v) \cdot \exp((t-t_0) \cdot v)} \\
	&= \exp(t_0 v) \cdot \diffd{}{s}\big|_{s=0} \exp(s v) \\
	&= \exp(t_0 v) \cdot v \\
	&= c_v(t_0) \cdot v = X_v \enbrace*{c_v(t_0)}_v
\end{align}

\begin{lemma}[label=lem:127]
	Sei $G$ eine Liegruppe.
	Die Integralkurve $c_v$ ist auf ganz $\mathbb{R}$ definiert und es gilt 
	\[
		c_v(t+s) = c_v(t) \cdot c_v(s)
	\]
\end{lemma}
\begin{beweis}
	Setze $\tilde{c}(t) := c_v(t_0) \cdot c_v(t)$.
	Es gilt 
	\begin{align}
		\tilde{c}'() = \enbrace*{\mathd L_{c_v(t_0)}}_{c_v(t)}  \cdot \Underbrace{\diffd{}{t} c_v(t)}{= X_v \enbrace*{c_v(t)} = \enbrace*{\mathd L_{c_v(t)}} \cdot v} = \enbrace*{\mathd L_{c_v(t_0) \cdot c_v(t)}}_e \cdot v = X_v \enbrace*{\tilde{c}(t)}
	\end{align}
	Somit ist $\tilde{c}(t)$ wieder eine Integralkurve von $X_v$.
	Beide Behauptungen folgen.
	Man nennt $(c_v(t))_{t \in \mathbb{R}}$ auch \Index{Einparameteruntergruppe} von $G$. 
\end{beweis}

\begin{definition}
	Sei $G$ eine Liegruppe.
	Dann nennt man 
	\[
		\exp \colon \mathfrak{g} \longrightarrow G ,\quad  v \longmapsto c_v(1)
	\]
	die \Index{Exponentialabbildung} von $G$. 
\end{definition}

\begin{lemma}
	Sei $G$ eine Liegruppe.
	Dann ist die Exponentialabbildung $\exp \colon \mathfrak{g} \to G$ differenzierbar und es gilt 
	\[
		\enbrace*{\mathd \exp}_0 = (\id)_{\Tmap_e G}
	\]
\end{lemma}
\begin{beweis}
	Wir betrachten das Vektorfeld
	\mapdef{S \colon G \times \mathfrak{g}}{\Tmap (G \times \mathfrak{g})}{(g,v)}{\enbrace*{X_v(g),0} \in \Tmap_g G \times \Tmap_v \mathfrak{g}}{}
	$S$ ist ein differenzierbares Vektorfeld mit globalem Fluss $\set*{\psi_t}_{t \in \mathbb{R}}$ gegeben durch 
	\[
		\psi_t(g,v) = \enbrace*{g \cdot \exp(t \cdot v),v}
	\]
	denn
	\begin{align}
		\diffd{}{t}\big| \psi_t (g,v) = \diffd{}{t} \enbrace*{g \cdot \exp(tv),v} = \enbrace*{\diffd{}{t} g \cdot \exp(tv),0} &= (\mathd L_g)_{c_{tv}(1)} \cdot \diffd{}{t} c_{tv}(1) \\
		&= \enbrace*{\mathd L_{g \cdot c_v(t_0)}}_e \cdot v \\
		&= X_v \enbrace*{g \cdot c_v(t_0)} = \exp(t_0 v)
	\end{align}
	Somit ist $\psi_t \colon G \times \mathfrak{g} \to G \times \mathfrak{g}$ differenzierbar für alle $t \in \mathbb{R}$.
	Insbesondere ist $\psi_1$ differenzierbar.
	Mit $\psi_1(e,\cdot ) = \enbrace*{\exp(\cdot ), \cdot }$ folgt, dass auch $\exp$ differenzierbar sein muss.
	
	Die zweite Behauptung ist klar.
\end{beweis}

\begin{satz}[label=satz:1210]
	Sei $G$ eine Liegruppe. Dann gilt
	\begin{enumerate}[1)]
		\item Eine Lieuntergruppe $H$ von $G$ ist genau dann eine eingebettete Untermannigfaltigkeit von $G$, wenn $H$ abgeschlossen ist.
		\item Eine \enquote{abstrakte} Untergruppe $H$ von $G$ ist eine Lieuntergruppe, falls $H$ abgeschlossen ist.
	\end{enumerate}
\end{satz}
\begin{beweis}
	\begin{enumerate}[1)]
		\item Sei $H \subset G$ ein eingebettete Lieuntergruppe. 
		Zu zeigen ist, dass $H$ abgeschlossen ist.
		Angenommen es existiert $(h_i)_{i \in \mathbb{N}} \subset H$ mit $h_i \to g \in G \setminus H$.
		Sei $U$ eine offene Umgebung von $e$, sodass $H \cap U$ eine zusammenhängende Untermannigfaltigkeit von $U$ ist.
		Eine solche Umgebung existiert, da $H$ eingebettet ist.
		Da $L_g$ ein Diffeomorphismus ist, ist $L_g(U)$ eine offene Umgebung von $g$.
		Somit existiert $i_0$ mit $h_i \in g(U)$ für alle $i \ge i_0$.
		Wegen $h_i \to g$ können wir $h_{i_0}^{-1} \cdot g \in U$ annehmen.
		Es gilt $h_{i_0}^{-1} \cdot g \notin H \cap U$, denn sonst folgt $g \in H$.
		Wegen $h_i \to g$ folgt
		\[
			\Underbracket{h_{i_0}^{-1}\cdot h_i}{\in H}  \xrightarrow{i \to \infty} h_{i_0}^{-1} \cdot g
		\]
		Dies ist -- in lokalen Koordinaten via einem Flachmacher -- ein Widerspruch und somit muss $H$ abgeschlossen sein.
		
		Sei nun $H$ eine abgeschlossene Lieuntergruppe von $G$.
		Zu zeigen ist, dass $H$ eine eingebettete Untermannigfaltigkeit ist, das heißt es existiert eine offene Umgebung $U$ von $e$ in $G$, sodass $H \cap U$ eine zusammenhängende Untermannigfaltigkeit von $U$ ist. 
		Annahme: Es existiert eine Umgebung $U$ von $e$, sowie eine Folge $(h_i)_{i \in \mathbb{N}}$ in $H$ mit
		\begin{itemize}
			\item $h_i \to e$
			\item $h_i \notin (U \cap H)_e$ (Zusammenhangskomponente der Eins)
		\end{itemize}
		Die Exponentialabbildung $\exp \Tmap_e G \to G$ ist nahe bei $e$ ein lokaler Diffeomorphismus (Umkehrsatz auf Mannigfaltigkeiten) mit 
		\[
			\enbrace*{\mathd \exp}_e = \id_{\Tmap_p M^n}
		\]
		Somit ist $h_i = \exp(v_i)$ für alle $i \in \mathbb{N}$ und $v_i \in \mathfrak{g}$ mit $v_i \to 0$ für $i \to \infty$.
		Wir wählen nun eine reelle Folge $(s_i)_{i \in \mathbb{N}}$ in $\mathbb{R}$ mit 
		\begin{equation}
			s_i \cdot v_i \to v \label{eq:1210:1} \tag{\#}
		\end{equation}
		längst einer Teilfolge.
		(Da $\mathfrak{g}$ ein Vektorraum ist, können wir das Skalarprodukt auf $\mathfrak{g}$ betrachten und erhalten eine Norm, setze dann $s_i := \frac{1}{\norm*{v_i}} $.)
		Sei $t \in \mathbb{R}$ \emph{fest} gewählt.
		\eqref{eq:1210:1} impliziert dann
		\begin{equation}
			t \cdot s_i \cdot v_i \grenzw{i \to \infty} t \cdot v \label{eq:1210:2} \tag{\#\#}
		\end{equation}
		Eine einfache Überlegung zeigt, dass man $t \cdot s_i \in \mathbb{Z}$ annehmen kann; 
		\[
			s_i \leadsto \tilde{s}_i := s_i + \text{\enquote{$\sfrac{1}{t}$}} 
		\]
		Beachte das \eqref{eq:1210:1} immernoch gilt.
		Aus $H \ni h_i = \exp(v_i)$ folgt 
		\[
			H \ni \exp(v_i)^{t \cdot s_i} = \exp \enbrace*{t \cdot s_i \cdot v_i} \xrightarrow[\text{\eqref{eq:1210:2}}]{i \to \infty} \exp(t \cdot v) \in H
		\]
		da $H$ abgeschlossen ist!
		Dies gilt für alle $t \in \mathbb{R}$.
		Somit $v \in \Tmap_e H$.
		Nächster Schritt: $\skal*{\cdot }{\cdot }$ sei ein Skalarprodukt auf $\Tmap_e G$.
		Dann gilt $\Tmap_e G = \mathfrak{h} \oplus_{\bot} \mathfrak{h}^{\bot}$.
		Mit dieser Zerlegung schreiben wir $v_i =x_i + y_i$.
		Die Abbildung 
		\mapdef{\Phi \colon \mathfrak{h} \oplus_{\bot} \mathfrak{h}^\bot}{G}{(x,y)}{\exp(x) \cdot \exp(y)}{}
		ist nahe $(0,0)$ ein lokaler Diffeomorphismus.
		Wir haben
		\[
			\Underbracket{\exp(v_i)}{\in H} = \Underbracket{\exp(\tilde{x}_i)}{\in H} \cdot \exp(\tilde{y}_i)
		\]
		$\leadsto e \leftarrow  \exp(\tilde{y}_i) = \exp(- \tilde{x}_i) \cdot \exp(v_i) \in H \setminus (H \cap U)_e$.
		Wie eben bewiesen, können wir nun eine Grenzrichtung 
		\[
			\frac{y_i}{\norm*{y_i}} \grenzw{i \to \infty}  y
		\]
		in $\mathfrak{h}$ konstruieren (mittels Teilfolge).
		Die Eigenschaften $y \in H$, $y \in H^\bot$ und $\norm*{y}=1$ widersprechen sich!
		Also existiert eine solche Folge nicht und $H$ ist eine eingebettete Untermannigfaltigkeit.
		\item Sei $H$ ein abgeschlossene \enquote{abstrakte} Untergruppe von $G$.
		Wir zeigen zunächst, dass 
		\[
			\mathfrak{h} := \set[\big]{v \in \Tmap_e G \given \exp(t \cdot v) \in H \,\forall t \in \mathbb{R}} \ni 0
		\]
		ein Unterraum ist.
		Annahme: Es gibt eine Folge $(h_i)_{i \in \mathbb{N}}$ in $H \setminus \set*{e}$ mit $h_i \to e$ (falls nicht: $e$ ist isoliert und $H$ ist eine diskrete Untergruppe von $G$).
		Wie in 1) zeigt man, dass die \enquote{Grenzrichtungen} $\lim_{i \to \infty} \frac{v_i}{\norm*{v_i}}$ in $\mathfrak{h}$ liegen.
		\begin{itemize}
			\item $v \in \mathfrak{h} \implies t \cdot v \in \mathfrak{h}$ für alle $t \in \mathbb{R}$
			\item $v_1 + v_2 \in \mathfrak{h} \implies v_1 + v_2 \in \mathfrak{h}$
		\end{itemize}
		Sei $(t_n)_{n \in \mathbb{N}}$ eine Nullfolge in $\mathbb{R}$.
		Dann gilt $t_n v_1, t_n v_2 \to 0$ für $n \to \infty$.
		Folglich gilt 
		\[
			\exp(t_n v_1), \exp(t_n v_2) \to e \text{ für }  n \to \infty
		\]
		Somit erhalten wir
		\[
			H \ni \exp(v_n) = \exp \enbrace*{t_n v_1} \exp(t_n v_2) \to e
		\]
		mit $v_n \to 0$.
		Nun gilt $\exp(t_n v_1) \cdot \exp(t_n v_2) = \exp \enbrace*{t_n(v_1 + v_2)} + \mathcal{O}(t_n^2)$.
		Somit
		\[
			\frac{v_n}{t_n} \grenzw{n \to \infty} v_1 +v_2 + \mathcal{O} \enbrace*{\frac{t_n^2}{t_n} \to 0 } 
		\]
		Damit ist $\mathfrak{h}$ ein Unterraum.
		Wie in 1) zeigt man nun, das $H$ ein Untermannigfaltigkeit ist.\qedhere
	\end{enumerate}
\end{beweis}

\begin{satz}[name={Ado},label=satz:1211]
	Jede endlichdimensionale Liealgebra $(V,\benbrace*{\cdot,\cdot }]$ ist isomorph zu einer Unteralgebra von $\mathfrak{gl}(n,\mathbb{R}) = \Tmap_e \GL(n,\mathbb{R})$.
\end{satz}

\begin{korollar}
	Zu jeder endlichdimensionalen Liealgebra $(V,\benbrace*{\cdot ,\cdot })$ existiert eine Liegruppe $G$ mit $V \cong \mathfrak{g}$
\end{korollar}
\begin{beweis}
	\autoref{satz:1211} und \autoref{satz:125}.
\end{beweis}

Bemerkung: Für $n \ge 3$ ist $\pi_1 \enbrace*{\SL(n,\mathbb{R})} = \mathbb{Z}_2$
Die universelle Überlagerung $\widetilde{\SL(n,\mathbb{R})}$ von $\SL(n,\mathbb{R})$ ist eine Liegruppe, welche in keine der Gruppen $\GL(N,\mathbb{R})$ eingebetten werden kann! Keine Matrixgruppe!

\begin{lemma}[label=lem:1213]
    Sei $G$ eine zusammenhängende Liegruppe und $U$ eine offene Umgebung von $e$ in $G$.
    Dann wird $G$ schon von $U$ erzeugt.
\end{lemma}
\begin{beweis}
    Sei $U \subset G$ offen mit $e \in U$.
    Dann existiert $V \subset U$ offen mit $e \in V$, sodass $V \cdot V^{-1} \subset U$, denn die Abbildungen $m \colon G \times G \to G$, $^{-1} \colon G \to G$ sind stetig.
    Wir setzen 
    \[
        H \coloneqq \bigcup_{n \in \mathbb{Z}} V^n
    \]
    $H$ ist eine offene Untergruppe von $G$.
    Zu zeigen bleibt, dass $H$ abgeschlossen ist.
    Die Abgeschlossenheit folgt wie im Beweis von \autoref{satz:125}.
\end{beweis}
% subsection 12 (end)

\section{Überlagerungen von Liegruppen} % (fold)
\label{sec:13}

\begin{definition}
    Sei $M^n$ eine $n$-dimensionale, zusammenhängende, differenzierbare Mannigfaltigkeit.
    Dann nennt man eine zusammenhängende, $n$-dimensionale, differenzierbare Mannigfaltigkeit $\tilde{M}^n$ \Index{Überlagerung} von $M^n$, falls eine surjektive differenzierbare Abbildung $\pi \colon \tilde{M} \to M$ existiert mit folgenden Eigenschaften:
    
    Für alle $p \in M$ existierte eine offene Umgebung $U$ von $p$ in $M^n$ mit folgenden Eigenschaften:
    \begin{enumerate}[1)]
        \item $\pi^{-1}(U) = \bigcup_{i \in I} \tilde{U}_i$ mit $\tilde{U}_i \subset \tilde{M}$ offen und $\tilde{U}_i \cap \tilde{U}_j = \emptyset$ für $i \neq j$.\marginnote{Jeder Punkt wird \enquote{schlicht überlagert}}
        \item $\pi\big|_{\tilde{U}_i} \colon \tilde{U}_i \to U$ ist ein Diffeomorphismus für alle $i \in I$.
    \end{enumerate}
    Die Abbildung $\pi \colon \tilde{M}^n \to M^n$ nennt man \Index{Überlagerungsabbildung} und $\abs*{I}$ die \Index{Blätterzahl} ($\abs*{I}= \infty$ ist möglich).
\end{definition}

Beispiele:
\begin{itemize}
    \item $M^n$ überlagert $M^n$.
    \item $S^3$ ist 2-fache Überlagerung von $\SO(3) = \mathbb{R}P^3$
    \item $S^n$ ist 2-fache Überlagerung von $\mathbb{R}P^n \coloneqq \set[\big]{\text{Menge der Geraden im }\mathbb{R}^{n+1}}$.
    \item Ist $G$ eine Gruppe, die frei und eigentlich diskontinuierlich auf $\tilde{M}^n$ operiert, so ist der Bahnenraum 
    \[
        M^n \coloneqq \set*{[p] \given p \in M^n} = \sfrac{\tilde{M}^n}{G}
    \]
    eine $n$-dimensionale, differenzierbare Mannigfaltigkeit und $\pi \colon \tilde{M}^n \to M^n$, $p \mapsto [p]$ eine Überlagerung
    \[
        [p] = [\tilde{p}] :\Longleftrightarrow \exists g \in G : p = g \cdot \tilde{p}
    \]
    Ist $G$ endlich: Dann ist $G$ frei und diskontinuierlich äquivalent zu: $g$ wirkt trivial auf einem Punkt $p$ genau dann, wenn $g$ auf allen $p$ trivial operiert.
    \item Für jede zusammenhängende Mannigfaltigkeit $M^n$ existiert eine bis auf Diffeomorphie einfach zusammenhängende Mannigfaltigkeit $\tilde{M}^n$, die $M^n$ überlagert (\Index{Universelle Überlagerung}).
    Ferner gilt $M^n = \sfrac{\tilde{M}^n}{\pi_1(M^n)}$.
\end{itemize}

\begin{lemma}[label=lem:132]
    Sei $G$ eine zusammenhängende Liegruppe.
    \begin{enumerate}[1)]
        \item Ist $\tilde{G}$ eine zusammenhängede Mannigfaltigkeit und $\pi \colon \tilde{G} \to G$ eine Überlagerungsabbildung, so besitzt $\tilde{G}$ genau eine Liegruppenstruktur, sodass $\pi$ ein Liegruppenhomomorphismus ist.
        \item Ein Liegruppenhomomorphismus $\Phi \colon \tilde{G} \to G$ mit $\tilde{G}$ zusammenhängend ist eine Überlagerung genau dann, wenn das Differential $(\mathd \Phi)_{\tilde{e}} \colon \Tmap_{\tilde{e}} \tilde{G} \to \Tmap_e G$ ein Isomorphismus ist.
    \end{enumerate}
\end{lemma}
\begin{beweis}
    \begin{enumerate}[(1)]
        \item Sei $\tilde{e} \in \pi^{-1}(e)$.
        Überlagerungstheorie impliziert (Stöcker-Zischer, Algebraische Topologie, Abschnitt 6.2)\todo{ref}, dass die stetigen Funktionen $m \colon G \times G \to G$, $^{-1} \colon G \to G$ entsprechende Lifts besitzen.
        Die Eindeutigkeit ist durch die Wahl von $\tilde{e}$ gegeben [Klar falls $\tilde{G}$ einfach zusammenhängend, sonst \enquote{kleines} Argument].
        
        Die Abbildungen $\tilde{m},\tilde{{ }^{-1}}$ sind differenzierbare Abbildungen, welche die Gruppenaxiome erfüllen: $\tilde{m} (\tilde{e}, \tilde{g}) = \tilde{g}$ für alle $\tilde{g} \in \tilde{G}$.
        Dies ergibt sich daraus, dass man Wege eindeutig liften kann.
        
        Die Abbildung $\pi \colon \tilde{G} \to G$ ist somit ein Liegruppenhomomorphismus.
        Die Eindeutigkeit der Liegruppenstruktur auf $\tilde{G}$, sodass $\pi \colon \tilde{G} \to G$ ein Liegruppenhomomorphismus, folgt wieder aus der Eindeutigkeit der Lifts.
        \item Die erste Implikation ist klar, denn Überlagerungsabbildungen sind lokale Diffeomorphismen.
        
        Für die Umkehrung sei $(\mathd \Phi)_{\tilde{e}} \colon \Tmap_{\tilde{e}} \tilde{G} \to \Tmap_e G$ ein Isomorphismus.
        Zu zeigen: Jedes $g \in G$ ist schlicht überlagert.
        Sei $\tilde{U}$ eine offene Umgebung von $\tilde{e}$ in $\tilde{G}$, sodass $\Phi|_{\tilde{U}} \colon \tilde{U} \to \Phi(\tilde{U})$ ein Diffeomorphismus ist (Umkehrsatz).
        Sei $\tilde{\Gamma} := \ker (\Phi)$.
        Dann gilt $\tilde{\Gamma} \cap \tilde{U}= \tilde{e}$.
        Sei nun $\tilde{V}$ eine offene Umgebung von $\tilde{e}$ mit $\tilde{V} \cdot \tilde{V}^{-1} \subset \tilde{U}$.
        \begin{itemize}
            \item  Seien $\tilde{\gamma}_1 \neq \tilde{\gamma}_2 \in \tilde{\Gamma}$ und $\tilde{v}_1, \tilde{v}_2 \in \tilde{V}$.
            Annahme:
            \[
                \tilde{\gamma}_1 \cdot \tilde{v}_1 = \tilde{\gamma}_2 \cdot \tilde{v}_2 \iff \tilde{\gamma}_2^{-1} \cdot \tilde{\gamma}_1 = \tilde{v}_2 \cdot \tilde{v}_1^{-1} \in \tilde{\Gamma} \cap \tilde{U}
            \]
            Mit obiger Eigenschaft des Schnitts folgt $\tilde{\gamma}_1 = \tilde{\gamma}_2$.
            Damit folgt $\tilde{\gamma}_1 \tilde{V} \cap \tilde{\gamma}_2 \tilde{V} = \emptyset$.
            \item Es gilt $\Phi^{-1}(V) = \bigcup_{\tilde{\gamma} \in \tilde{\Gamma}} \tilde{\gamma}(v)$:
            
            Sei $\tilde{v} \in  \tilde{V}$ und $\tilde{g} \in \tilde{G}$ mit $\Phi(\tilde{v}) = \Phi(\tilde{g})$.
            Dies ist äquivalent zu
            \[
                \Phi(\tilde{g}^{-1} \tilde{v}) = \Phi(\tilde{g}^{-1}) \Phi(\tilde{v}) = e
            \]
            Also ist $\tilde{g}^{-1} \tilde{v} = \tilde{\gamma} \in \tilde{\Gamma}$ und somit $\tilde{v} = \tilde{\gamma} \cdot \tilde{g}$.
            \item Die Abbildung $\Phi|_{\tilde{\gamma} \tilde{V}} \colon \tilde{\gamma}(\tilde{V}) \to V$ ist ein Diffeomorphismus:
            
            \[
                \Phi(\tilde{\gamma} \tilde{v}) = \Phi(\tilde{\gamma}) \cdot \Phi(\tilde{v}) = \Phi \enbrace*{\tilde{\gamma}^{-1} \tilde{\gamma} \tilde{v}} = \enbrace*{\Phi \circ L_{\tilde{\gamma}^{-1}}} (\tilde{\gamma} \tilde{v})
            \]
            \item Alle $g \in G$ sind schlicht überlagert.
            \item Die Abbildung $\Phi$ ist surjektiv nach \autoref{lem:1213}, denn eine offene Umgebung von $e$ liegt im Bild und $\Phi$ ist ein Homomorphismus.\qedhere
        \end{itemize}
    \end{enumerate}
\end{beweis}

\begin{korollar}[label=lem:133]
    Seien $\tilde{G}$, $G$ zusammenhängende Liegruppen. 
    \begin{enumerate}[1)]
        \item Ist $\Phi \colon \tilde{G} \to G$ eine Überlagerung von Liegruppen, so ist $\tilde{\Gamma} = \ker \Phi$ eine diskrete Untergruppe von $Z(G)$.\marginnote{Zentrum von $\tilde{G}$}
        \item Ist $\tilde{\Gamma}$ eine diskrete Untergruppe von $Z(\tilde{G})$, so ist $G \coloneqq \sfrac{\tilde{G}}{\tilde{\Gamma}}$ eine Liegruppe und $\Phi \colon \tilde{G} \to G$, $\tilde{g} \mapsto [\tilde{g}]$ eine Überlagerung.
    \end{enumerate}
\end{korollar}
\begin{beweis}
    \begin{enumerate}[1)]
        \item Sei $\tilde{U}$ eine offene Umgebung von $\tilde{e}$, sodass $\Phi|_{\tilde{U}} \colon \tilde{U} \to \Phi(\tilde{U}) =: U$ ein Diffeomorphismus ist.
        Dann folgt $\tilde{\Gamma} \cap \tilde{U} = \set*{\tilde{e}}$.
        Somit ist $\tilde{\Gamma}$ eine diskrete Untergruppe von $\tilde{G}$.
        Da $\tilde{\Gamma}$ der Kern eine Homomorphismus ist, ist $\tilde{\Gamma}$ ein Normalteiler in $\tilde{G}$.
        Sei $\tilde{g} \in \tilde{G}$ und $\tilde{g}(t) \colon [0,1] \to \tilde{G}$ ein stetiger Weg mit $\tilde{g}(0)=\tilde{e}$, $\tilde{g}(1)=\tilde{g}$.
        Dann gilt
        \[
            \tilde{g}(t) \tilde{\gamma} \tilde{g}^{-1}(t) \in \tilde{\Gamma}
        \]
        Es folgt -- da $\tilde{\Gamma}$ diskret ist -- $\tilde{g}(t) \tilde{\gamma} (\tilde{g}(t))^{-1} \equiv \tilde{\gamma}$.
        Es folgt 
        \[
            \tilde{g} \tilde{\gamma} \tilde{\gamma} \tilde{g} \forall \tilde{g} \in \tilde{G} \iff \tilde{\gamma} \in Z(\tilde{G})
        \]
        \item \emph{Übung!}
    \end{enumerate}
\end{beweis}

Ein kleines Beispiel:
\[
    Z \enbrace*{\SO(n)} = \begin{cases}
        \set*{\id_n} &\text{ falls } n \text{ ungerade} \\
        \set*{\pm \id_n} &\text{ falls $n$ gerade}
    \end{cases}
\]
Für $n \ge 2$
\[
    Z \enbrace*{\SU(n)} = \mathbb{Z}_n = \set*{e^{2 \pi i \sfrac{k}{n}} \cdot \id_n \given k=0,1,\ldots ,n-1 }
\]
Es gilt $\pi_1(\SO(n))= \mathbb{Z}_2$.
Die universelle Überlagerung ist $\mathrm{Spin}(n)$.
Für $n\ge 2$ $\pi \enbrace*{\SU(n)} = \set*{e}$ und $\SU(2) = S^3$.

\begin{lemma}[label=lem:134]
	Seien $H,G$ zusammenhängende Liegruppen.
	Ist $H$ einfach zusammenhängend und $\psi \colon \Tmap_e H \to \Tmap_e G$ ein Liealgebrenhomomorphismus, so existiert ein eindeutiger Liegruppenhomomorphismus $\Phi \colon H \to G$ mit $(\mathd \Phi)_e = \psi$.
\end{lemma}
\begin{beweis}
	Da der Graph
	\[
		\Graph(\psi) = \set[\big]{\enbrace*{v,\psi(v)} \given v \in \Tmap_e H} \subset \Tmap_e H \oplus \Tmap_e G
	\]
	eine Unteralgebra von $\Tmap_e H \oplus \Tmap_e G$ ist, existiert nach \autoref{satz:125} eine eindeutig bestimmte Lieuntergruppe $A$ von $H \times G$ mit $\Tmap_e A = \Graph(\psi)$.
	Wir bezeichnen mit $\pi_1 \colon A \to H$ und $\pi_2 \colon A \to G$ die Projektionen auf den ersten bzw. zweiten Faktor von $H \times G$.
	Die Abbildungen $\pi_1$ und $\pi_2$ sind Liegruppenhomomorphismen.
	Ferner ist 
	\mapdef{(\mathd \pi_1)_e \colon \Tmap_e A}{\Tmap_e H}{(v,\psi(v))}{v}{}
	ein Isomorphismus.
	Nach \autoref{lem:132} 2) überlagert $A$ somit $H$.
	Da $H$ einfach zusammenhängend ist, folgt $A \cong H$.
	Wir erhalten somit einen Homomorphismus 
	\mapdef{\Phi \colon H = A}{G}{h}{\pi_2(h)}{}
	mit $(\mathd \Phi)_e = (\mathd \pi_2)_e = \psi$.
\end{beweis}

Ist in obiger Situation $H$ \emph{nicht} einfach zusammenhängend, so erhält man einen Homomorphismus $\tilde{\Phi} \colon \tilde{H} \to G$, wobei $\tilde{H}$ eine universelle Überlagerung von $H$ ist.
Dies liefert genau dann einen Homomorphismus $\Phi \colon H \to G$, wenn $\ker \pi \subset \ker \tilde{\Phi}$ ist, wobei $\pi \colon \tilde{H} \to H$ die Überlagerungsabbildung ist.

Dies ist aber nicht immer richtig:
So kann man zum Beispiel den durch die Identität gegeben Liealgebrenisomorphismus $\psi \colon \mathfrak{h} \to \mathfrak{h}$ betrachten.
Dann ist die induzierte Abbildung $\tilde{\Phi} \colon \tilde{H} \to \tilde{H}$ injektiv, aber $\pi \tilde{H} \to H$ muss nicht injektiv sein.
Ein konkretes Beispiel wäre die universelle Überlagerung von $\SL(n,\mathbb{R})$.

\begin{satz}
	Es gilt
	\begin{enumerate}[1)]
		\item Zwei zusammenhängende, einfach zusammenhängenden Liegruppen mit isomorphen Liealgebren sind isomorph.
		\item Für jede endlich dimensionale Liealgebra $(V, \benbrace*{\cdot,\cdot })$ existiert genau eine zusammenhängende, einfach zusammenhängende Liegruppe $\tilde{G}$ mit Liealgebra $\tilde{\mathfrak{g}}=V$.
		\item Jede zusammenhängende Liegruppe $G$ ist isomorph zu $\sfrac{\tilde{G}}{\tilde{\Gamma}}$, wobei $\tilde{G}$ eine zusammenhängende, einfach zusammenhängende Liegruppe und $\tilde{\Gamma} \subset Z(\tilde{G})$ eine diskrete Untergruppe ist.
	\end{enumerate}
\end{satz}
\begin{beweis}
	\begin{enumerate}[1)]
		\item Sei $\psi \colon \Tmap_e H \to \Tmap_e G$ ein Liealgebrenisomorphismus, $H,G$ zusammenhängende und einfach zusammenhängende Liegruppen.
		Nach \autoref{lem:134} existieren Homomorphismen $\Phi_1 \colon H \to G$ und $\Phi_2 \colon G \to H$ mit
		\[
			(\mathd \Phi_1)_e = \psi \qquad (\mathd \Phi_2)_e = \psi^{-1}
		\]
		Somit ist $\Phi_2 \circ  \Phi_1 \colon H \to G$ ein Liegruppenhomomorphismus mit $(\mathd \Phi_2 \circ \Phi_1)_e = \id_{\Tmap_e H}$.
		Nach \autoref{kor:126} gilt $\Phi_2 \circ \Phi_1= \id_H$ und somit $\Phi_2 = \Phi_1^{-1}$.
		\item Nach dem Satz von Ado existiert eine Liegruppe $G$ mit Liealgebra $\mathfrak{g}=V$.
		Die universelle Überlagerung $\tilde{G}$ von $G$ ist die gesuchte Liegruppe.
		Eindeutigkeit folgt aus 1).
		\item \autoref{lem:133}\qedhere
	\end{enumerate}
\end{beweis}

Beispiel: 
\[
	\SL(2,\mathbb{R}) = \set*{\begin{pmatrix}
		a & b \\ c & d
	\end{pmatrix} \given ad - bc =1}
\]
ist eine 3-dimensionale Untergruppe von $\GL(n,\mathbb{R})$.
Es gilt $\SL(2,\mathbb{R}) = \SO(2) \times \mathbb{R}^2$.
Somit ist $\pi_1 \enbrace*{\SL(2,\mathbb{R})} = \pi_1 (\SO(2)) = \mathbb{Z}$.
Es sei $M^3 := D^2 \times \mathbb{R} \cong \mathbb{R}^3$ (diffeomorph) mit
\[
	D^2 := \set*{z \in \mathbb{C} \given \abs*{z} <1} = \set*{\begin{pmatrix}
		x \\ y
	\end{pmatrix} \in \mathbb{R}^2 \given x^2 +y^2 <1}
\]
Wir führen folgende Multiplikation auf $M^3$ ein:
\[
	(\alpha,\varphi) \cdot (\beta, \psi) := \enbrace*{ \frac{\beta + \alpha \cdot e^{-i \psi}}{1 + \alpha \tilde{\beta} e^{-i \psi}}, \varphi + \psi + 2 \mathrm{Arg} \enbrace*{1 + \alpha \overline{\beta} e^{-i \psi}} }
\]
Es gilt zum Beispiel $(\alpha,\varphi) \cdot (0,0) = \enbrace*{\alpha,\varphi} = (0,0) \cdot (\alpha,\varphi)$.
Das heißt das neutrale Element ist $e=(0,0)$.
Ferner gilt
\[
	(\alpha,\varphi) \cdot \enbrace*{- \alpha e^{i \varphi}, -\varphi} = (0,0) = e
\]
Somit ist $M^3$ eine einfach zusammenhängende Liegruppe (!).
Man rechnet leicht nach, dass 
\begin{align}
	E_1(\alpha,\varphi) &= \enbrace*{\frac{1}{2}  (1- \alpha^2), \Im \alpha} = \enbrace*{\frac{1}{2} (1 -x^2 +y^2), -xy,y} \\
	E_2(\alpha,\varphi) &= \enbrace*{\frac{1}{2} i (1+ \alpha^2), - \Re(\alpha) } = \enbrace*{-xy, \frac{1}{2} \enbrace*{1+x^2 -y^2}, -x} \\
	E_3(\alpha,\varphi) &= \enbrace*{-i \alpha,1} = (y,-x,1)
\end{align}
linksinvariante Vektorfelder mit
\[
	E_1(0,0) = \enbrace*{\sfrac{1}{2},0,0 } \qquad E_2(0,0) = \enbrace*{0, \sfrac{1}{2}, 0 } \qquad E_3(0,0) = (0,0,1)
\]
Dies erhält man zum Beispiel durch
\[
	E_1(\alpha,\varphi) = \mathd L_{(\alpha,\varphi)} \cdot \enbrace*{\sfrac{1}{2},0,0} = \diffd{}{t}\big|_{t=0} (\alpha,\varphi) \cdot \enbrace*{\sfrac{1}{2},0,0} = \ldots 
\]
Es gilt
\[
	\benbrace*{E_1,E_2} = - E_3 \quad \benbrace*{E_3,E_1} = E_2 \quad \benbrace*{E_2,E_3} =E_1
\]
Auf $\SL(2,\mathbb{R})$
\[
	\overline{E}_1 := \frac{1}{2} \begin{pmatrix}
		1 & 0 \\ 0 & -1
	\end{pmatrix} \quad 
	\overline{E}_2 := \frac{1}{2}  \begin{pmatrix}
		0 & 1 \\ 1 & 0
	\end{pmatrix} \quad 
	\overline{E}_3 := \frac{1}{2} \begin{pmatrix}
		0 & -1 \\ 1 & 0
	\end{pmatrix} 
\]
% section 13 (end)

\section{Die Exponentialabbildung} % (fold)
\label{sec:14}

\begin{definition}[{name=[{Einparameteruntergruppe}]}]
	Sei $G$ eine Liegruppe.
	Dann nennt man einen Homomorphismus $\Phi \colon (\mathbb{R},+) \to G$ eine \Index{Einparameteruntergruppe} (von $G$).
\end{definition}

Beispiel: $G = \GL(n;\mathbb{R})$, $v \in \Tmap_e \GL(n,\mathbb{R}) = \Mat(n,\mathbb{R})$.
\[
	\Phi_v \colon \mathbb{R} \longrightarrow G \qquad t \longmapsto \exp(t \cdot v) = \sum_{k=0}^{\infty}  \frac{(tv)^k}{k!} 
\]
Es gilt $\Phi_v(s+t) = \Phi_v(s) \cdot \Phi_v(t) = \Phi_v(t) \cdot \Phi_v(s)$ für alle $s,t \in \mathbb{R}$, wie man sich leicht überlegt.

Zurück zum allgemeinen Fall: Sei $v \in \Tmap_e G \setminus \set*{0}$.
Wegen $\Tmap_0 \mathbb{R} \cong \mathbb{R}$ ist
\[
	\psi_v \colon \mathbb{R} \to \Tmap_e G \quad t \mapsto t \cdot v
\]
ein Liealgebrenhomomorphismus. 
Da $\mathbb{R}$ einfach zusammenhängend ist, existiert nach \autoref{lem:134} ein eindeutig bestimmter Homomorphismus
\(
	\Phi_v \colon \mathbb{R} \to G
\)
mit $\Phi_v'(0)= (\mathd \Phi_v)_0 = \psi_v$.

\begin{lemma}[label=lem:142]
	Sei $G$ eine Liegruppe mit Exponentialabbildung $\exp \colon \Tmap_e G \to G$.
	Dann gilt $\Phi_v(1) = \exp(v)$ für alle $v \in \Tmap_e G$.
\end{lemma}
\begin{beweis}
	Wir hatten in \autoref{lem:127} gezeigt, dass $c_v(s +t) = c_v(s) \cdot c_v(t)$ für alle $t,s \in \mathbb{R}$ gilt, wobei $c_v$ die Integralkurve von $X_v$ ist mit $X_v(e)=v$.
	Somit gilt
	\[
		\exp \enbrace*{(s+t) v} = c_{(s+t)v}(1) = c_v(s+t) = c_v(s) \cdot c_v(t) = c_{sv}(1) \cdot c_{tv}(1) = \exp(sv) \cdot \exp(tv)
	\]
	Damit ist $t \mapsto \exp(tv)$ eine Einparametergruppe mit $\diffd{}{t}\big|_{t=0} \exp(tv)=v$.
	Die Eindeutigkeit impliziert $\Phi_v(t)=\exp(tv)$.
\end{beweis}

\begin{lemma}[label=lem:143]
	Es gilt
	\begin{enumerate}[(1),itemsep=0pt]
		\item Für alle $v \in \Tmap_eG$ ist $\Phi_v(t)=\exp(tv)$ eine Einparamtergruppe von $G$ mit $\Phi'_v(0)=v$.
		\item Sei $g \in G$ und $v \in \Tmap_eG$.
		Dann ist $L_g(\Phi_v(t))$ Integralkurve von $X_v$.
		\item Ist $\Phi \colon H \to G$ ein Homomorphismus, so gilt für alle $w \in \Tmap_eW$ für alle $t \in \mathbb{R}$
		\[
			\Phi \enbrace*{\exp_H(tw)} = \exp_G \enbrace*{t \enbrace*{\mathd \Phi}_e \cdot w}
		\]
		\item Ist $H$ eine Lieuntergruppe von $G$, so gilt
		\[
			\Tmap_e H = \set*{v \in \Tmap_e G \given \exp_G(tv) \in H \enspace \forall t \in \mathbb{R}}
		\]
	\end{enumerate}
\end{lemma}
\begin{beweis}
	\begin{enumerate}[(1),itemsep=0pt]
		\item siehe \autoref{lem:142}
		\item Klar, denn $\exp(tv)=c_v(t)$.
		\item Homomorphismen bilden Einparametergruppen auf Einparametergruppen ab:
		\[
			\Phi \enbrace*{h(s+t)} = \Phi \enbrace*{h(s) \cdot h(t)} = \Phi \enbrace*{h(s)} \Phi \enbrace*{h(t)}
		\]
		mit $\diffd{}{t}\big|_{t=0} (\Phi \circ h)(t) = \enbrace*{\mathd\Phi}_e \cdot h'(0)$.
		Aus der Eindeutigkeit folgt die Behauptung.
		\item Sei $i \colon H \to G$ die Inklusion.
		Dann ist $\exp_H(tw) = i \enbrace*{\exp_H(tw)} = \exp_G \enbrace*{(\mathd i)_e \cdot w \cdot t} = \exp_G(tw)$.\qedhere
	\end{enumerate}
\end{beweis}

Beispiel: Es gilt $\Tmap_e\SO(n) = \set*{V \in \Mat_n(\mathbb{R}) \given V^T =-V}$
\[
	\benbrace*{V_1,V_2}^T = \enbrace*{V_1 V_2 - V_2 V_1}^T = V_2 V_1 - V_1 V_2 = -\benbrace*{V_1,V_2}
\]
Aus (4) folgt $\exp(tv) \in \SO(n)$ für alle $t \in \mathbb{R}$. Anders: $\exp_H = \exp_G\big|_{\Tmap_e H} \colon \Tmap_e H \to H$.

Bemerkung: 
\begin{enumerate}[1)]
	\item ist $H$ eine Lieuntergruppe von $G$, so gilt $\exp_H  = \exp_B\big|_{\Tmap_e H} \colon \Tmap_e H \to H$
	\item Die Exponentialabbildung $\exp \colon \Tmap_e G \to G$ ist nicht immer surjektiv; richtig für kompakte zusammenhängende Liegruppen.
	Es gilt $\im \enbrace*{\exp (\Tmap_e G)} \supset U \supset \set*{e}$ für eine offene Umgebung von $e$.
	\item Es gilt $\exp(x)^{-1} = \exp(-x)$.
	\item Der Fluss eines linksinvarianten Vektorfeldes $X_v$ ist gegeben durch
	\[
		\enbrace*{R_{\exp(tv)}}_{t \in \mathbb{R}} \qquad X_v(g) = \enbrace*{\mathd L_g}_e \cdot v
	\]
\end{enumerate}
% section 14 (end)

\section{Die adjungierte Darstellung} % (fold)
\label{sec:15}
Sei $G$ eine Liegruppe und $g \in G$ fest. Betrachte
\mapdef{i_g \colon G}{G}{\tilde{g}}{m \enbrace*{m \enbrace*{g,\tilde{g}}, I(g)}= g \cdot \tilde{g} \cdot g^{-1}}{}
wobei $m(g,\tilde{g})= g \cdot \tilde{g}$, $I(g)=g^{-1}$. 
Somit gilt
\[
	i_g = L_g \circ R_{g^{-1}} = R_{g^{-1}} \circ L_g
\]
Wir hatten bereits gesehen, dass $i_g \colon G \to G$ ein Liegruppenhomomorphismus ist.
Folglich ist 
\[
	\Ad(g) := \enbrace*{\mathd i_g}_e \colon \Tmap_e G \longrightarrow \Tmap_eG = \mathfrak{g}
\]
ein Liealgebrenhomomorphismus.
Wir hatten auch gezeigt, dass aus $i(g_1 \cdot g_2) = i(g_1) \cdot i(g_2)$ die Identität 
\[
	\Ad(g_1 \cdot g_2) = \Ad(g_1) \circ \Ad(g_2)
\]
folgt.

\begin{definition}[{name=[adjungierte Darstellung]}]
	Sei $G$ eine Liegruppe,
	Dann nennt man 
	\[
		\Ad \colon G \to \GL(\mathfrak{g}) \qquad g \longmapsto \Ad(g)
	\]
	die \Index{adjungierte Darstellung} von $G$.
\end{definition}

Ist zum Beispiel $G$ abelsch, so ist $i_g = \id_G$ und es folgt ${\Ad(g)} = \id_{\mathfrak{g}}$.
Ist $G=S^3$, so wurde in den Übungen die adjungierte Darstellung von $G$ bestimmt:
\[
	\Tmap_e S^3 = \im \mathbb{H} = \Span_\mathbb{R} (i,j,k) =: \mathbb{R}^3
\]
Es gilt $\Ad(g)(x) = g \cdot x \cdot g^{-1} \in \mathbb{R}^3$.
Es wurde auch $\Im \enbrace*{\Ad(S^3)} = \S0(3)$ gezeigt.

\begin{definition}
	Sei $G$ eine Liegruppe mit Liealgebra $\mathfrak{g}$.
	Dann bezeichnen wir mit 
	\[
		\ad_v \colon \mathfrak{g} \to \mathfrak{g} \qquad w \longmapsto \benbrace*{v,w}
	\]
	\emph{hier fehlt noch ein Satz irgendwie}
\end{definition}

Die Jacobi-Identität 
\[
	\benbrace*{X,\benbrace*{Y,Z}} + \benbrace*{Z, \benbrace*{X,Y}} + \benbrace*{Y, \benbrace*{Z,X}} =0
\]
für alle $X,Y,Z$ kann nun folgendermaßen beschrieben werden:\marginnote{Kommutator}
\[
	\ad_{\benbrace*{X,Y}} = \ad_X \circ \ad_Y - \ad_Y \circ \ad_X = \benbrace*{\ad_X, \ad_Y}_0
\]
Folglich ist $\ad \colon \enbrace*{\mathfrak{g}, \benbrace*{\cdot ,\cdot }} \to \enbrace*{\mathfrak{gl}(\mathfrak{g}), \benbrace*{\cdot ,\cdot }_0}$
ein Liealgebrenhomomorphismus.

\begin{lemma}
	Sei $G$ eine Liegruppe.
	Dann gilt
	\begin{enumerate}[(1)]
		\item $\enbrace*{\mathd \Ad}_{\id_{\mathfrak{g}}} = \ad$
		\item $i_g \enbrace*{\exp(v)} = \exp \enbrace*{\Ad(g)v}$
		\item $\Ad \enbrace*{\exp(v)} = e^{\ad_v} = \sum_{k=0}^{\infty} \frac{(\ad_v)^k}{k!} \colon \mathfrak{g} \to \mathfrak{g}$
		\item Ist $G$ zusammenhängend, so gilt $Z(G) = \ker (\Ad)$
	\end{enumerate}
\end{lemma}
\begin{beweis}
	\begin{enumerate}[(1)]
		\item Erinnerung: Ist $M^n$ eine differenzierbare Mannigfaltigkeit, $p \in M^n$ ein Punkt, $X,Y$ glatte Vektorfelder und $\enbrace*{\psi_t}_{\abs*{t}<\varepsilon}$ lokaler Fluss von $X$ in einer Umgebung von $p$.
		Dann gilt
		\[
			\benbrace*{X,Y}(p) = \diffd{}{t}\Big|_{t=0} \enbrace*{\mathd \psi_{-t}}_{\psi_t(p)} Y \enbrace*{\psi_t(p)}
		\]
		Es folgt für alle $w \in \mathfrak{g}$.
		\begin{align}
			\enbrace*{\mathd \enbrace*{\Ad(g)}_{\id_{\mathfrak{g}}} \cdot v} (w) = \diffd{}{t}\Big|_{t=0} \Ad \enbrace*{\exp(tv)}(w) &= \diffd{}{t}\Big|_{t=0} \enbrace*{\mathd R_{\exp(-tv)}}_{\exp{tv}} \cdot  \enbrace*{\mathd L_{\exp(tv)}}_{e} \cdot w \\
			&= \diffd{}{t}\Big|_{t=0} \enbrace*{\mathd\enbrace*{R_{\exp(-tv)}}}_{\exp{tv}} X_w \enbrace*{\exp(tv)} \\
			&= \benbrace*{X_v,X_w}(e)  = \benbrace*{v,w}_{\mathfrak{g}} \\
			&= \ad_v(w)
		\end{align}
		\item Folgt direkt aus \autoref{lem:143}.
		\item Folgt direkt aus \autoref{lem:143}.
		\item Zur Inklusion $Z(G) \subseteq ker (\Ad)$: Ist $g \in Z(g)$, so ist $i_g = {\id_G} $ und es folgt $\Ad(g)=\id$.
		
		Es gelte andersrum $\Ad(g) = \id_{\mathfrak{g}}$.
		Dann sind $i_g, \id_G \colon G \to G$ Liegruppenhomomorphismen mit $\enbrace*{\mathd i_g}_e = \mathd(\id_G)$.
		Da $G$ zusammenhängend ist, folgt aus \autoref{kor:126} $i_g=\id_G$ und es folgt $g \in Z(G)$.
	\end{enumerate}
\end{beweis}

Beispiel: Betrachte $G= \GL(n,\mathbb{R})$.
Dann gilt
\[
	\begin{array}{ll}
		i_A(\tilde{A}) = A \cdot \tilde{A} \cdot A^{-1} \quad & \tilde{A} \in \GL(n,\mathbb{R}) \\
		\Ad(A)(B) = A \cdot B \cdot A^{-1} & B \in \Mat(n,\mathbb{R}) \\
		\ad_v(W) = V \cdot W + W (-V) = \benbrace*{V,W}_0
	\end{array}
\]

\begin{lemma}
	Sei $G$ eine zusammenhängende Liegruppe. Dann gilt
	\begin{enumerate}[1)]
		\item $G$ ist abelsch genau dann, wenn $\Tmap_e G = \mathfrak{g}$ abelsch ist.
		\item $G$ ist genau dann abelsch, wenn $G \cong T^k \times \mathbb{R}^{n-k}$.
		\item Ist $H$ eine zusammenhängende Untergruppe von $G$, so ist $H$ Normalteiler genau dann, wenn $\mathfrak{h} = \Tmap_e H$ Ideal von $\mathfrak{g}$ ist, das heißt wenn $\benbrace*{\mathfrak{h},\mathfrak{h}} \subset \mathfrak{h}$ gilt.
		\item Das Zentrum $Z(G)$ ist Lieuntergruppe von $G$ mit 
		\[
			Z(\mathfrak{g}) = \Tmap_e Z(G) = \set*{v \in \mathfrak{g} \given \ad_v =0}
		\]
		\item Ist $H$ eine Lieuntergruppe von $G$, so ist die Einschränkung der adjungierten Darstellung von $G$ auf $H$ gerade die adjungierte Darstellung von $H$.
	\end{enumerate}
\end{lemma}
\begin{beweis}
	\begin{enumerate}[1)]
		\item Sei $G$ abelsch. Dann gilt $i_g=\id_G$, also $\Ad(g) = \id_{\mathfrak{g}}$ für alle $g \in G$ und es folgt $\ad_v=0$. Damit ist auch die Lieklammer trivial und $\mathfrak{g}$ ist abelsch.
		
		Sei andersrum $\mathfrak{g}$ abelsch: Dann ist $\ad_v=0$ für alle $v \in \mathfrak{g}$. Es gilt $\id_{\mathfrak{g}} = e^{\ad_v} = \Ad \enbrace*{\exp(v)}$
		Es folgt $\Ad(g) = \id_{\mathfrak{g}} $ für alle $g \in G$.
		\[
			i_g \enbrace*{\exp(v)} =  g \exp(v) g^{-1} = \exp \enbrace*{\Ad(g)v} = \exp(v)
		\]
		Somit ist $G$ abelsch.
		\item Da $G$ abelsch ist, ist $\mathfrak{g}$ abelsch und somit ist $\exp \colon (\Tmap_eG,+) \to (G,\cdot )$ ein Liegruppenhomomorphismus (wenn der Kommutator verschwindet gilt die Funktionalgleichung).
		Somit ist $\exp \colon \Tmap_e G \to G$ eine Überlagerung und $\Gamma := \exp^{-1}(e)$ eine diskrete Untergruppe von $(\Tmap_e G,+)$.\marginnote{Gitter}
		Somit existieren $v_1, \ldots, v_k \in \mathfrak{g}$ mit
		\[
			\Gamma = \set*{\sum\nolimits_{i=1}^{k} \lambda_i v_i \given \lambda_i \in \mathbb{Z}} \cong \mathbb{Z}^k
		\]
		Sei $V := \Span_\mathbb{R}(v_1,\ldots ,v_k)$ und $V^\bot$ ein Komplement von $V$ in $\mathfrak{g}$.
		Es gilt also $\mathfrak{g} = V \oplus V^\bot$ (beide Grenzfälle sind möglich).
		Es ist nun \enquote{klar}, dass 
		\[
			\im \enbrace*{\exp(V)} = T^k
		\]
		und $\im \enbrace*{\exp(V^\bot)} = \mathbb{R}^{n-k}$ gilt. 
		Die Behauptung folgt.
		\item Sei $H$ ein zusammenhängender Normalteiler von $G$, das heißt es gilt $g H g^{-1} \subset H$ für alle $g \in G$.
		Es folgt $i_g(H) \subset H$ und damit auch $\Ad(g)(\Tmap_eH) \subset \Tmap_eH$ für alle $g \in G$.
		Damit gilt auch $\ad_v(\Tmap_e H) \subset \Tmap_e H$, also $\benbrace*{\mathfrak{g},\mathfrak{h}}\subset \mathfrak{h}$.
		
		Ist umgekehrt $\mathfrak{h} = \Tmap_e H$ ein Ideal von $\mathfrak{g}$, so gilt $\ad_v(\mathfrak{h}) \subset \mathfrak{h}$ für alle $v \in \mathfrak{g}$.
		Dann gilt auch $\Ad \enbrace*{\exp (v)} = e^{\ad_v}(\mathfrak{h}) \subset \mathfrak{h}$.
		Damit gilt auch $\Ad(g)(\mathfrak{h}) \subset \mathfrak{h}$ für alle $g \in G$.
		\begin{align}
			H \ni \exp_G \enbrace*{\Ad(g)(w)} = i_g \enbrace*{\exp(w)}= g \cdot \exp(w) \cdot g^{-1} 
		\end{align}
		für alle $g \in G$ und $w \in \mathfrak{g}$.
		Damit ist $H$ ein Normalteiler.
		\item \emph{Übung!}
		\item \emph{Übung!}\qedhere
	\end{enumerate}
\end{beweis}

Frage: Wann sind Normalteiler abgeschlossene Untergruppen?
Für $r$ irrational ist
\[
	S^1_r = \set*{\enbrace*{e^{i \varphi}, e^{i r \varphi}} \in T^2}
\]
ist Normalteiler, aber $\overline{S^1_r}=T^2$.
% section 15 (end)

\section{Automorphismen von Liegruppen} % (fold)
\label{sec:16}

\begin{definition}
	Sei $G$ eine Liegruppe mit Liealgebra $\mathfrak{g}$.
	\begin{enumerate}[(i)]
		\item Dann nennt man einen linearen Isomorphismus $A \colon \mathfrak{g} \to \mathfrak{g}$ einen Liealgebra-Isomorphismus, falls 
	\[
		A \benbrace*{X,Y} = \benbrace*{AX,AY}
	\]
	für alle $X,Y \in \mathfrak{g}$ gilt.
	Diese Menge der Automorphismen bezeichnen wir mit $\Aut(\mathfrak{g}) \subset \GL(\mathfrak{g})$.
	\item Eine lineare Abbildung $A \colon \mathfrak{g} \to \mathfrak{g}$ nennt man \Index{Derivation}, falls 
	\[
		A \benbrace*{X,Y} = \benbrace*{AX,Y} + \benbrace*{X,AY}
	\] 
	für alle $X,Y \in \mathfrak{g}$ gilt. Bezeichnung: $\Der(\mathfrak{g}) \subset \End(\mathfrak{g})$.
	\end{enumerate}	
	Automorphismen sind Untergruppe von $\GL(\mathfrak{g})$
\end{definition}

% section 16 (end)
% chapter 1 (end)

\cleardoubleoddemptypage
\pagenumbering{Alph}
\setcounter{page}{1}
\cleardoubleoddemptypage
\appendix

\chapter{Anhang} % (fold)
\label{sec:anhang}
%!TEX root = ana_top_geo.tex

\subsection{Ausführlicher Beweis zu \cref{lem:kpt-schnitte}} % (fold)
\label{sub:kpt-schnitte}
Sei $X$ ein Hausdorffraum. Dann ist $X$ genau dann kompakt, wenn gilt: Hat eine Familie $\mathcal{A}$ von abgeschlossenen Teilmengen von $X$ die endliche 
Durchschnittseigenschaft, so gilt 
\[
	\bigcap_{A \in \mathcal{A}} A \not= \emptyset.
\]
\begin{beweis}
	Für die erste Implikation sei $X$ kompakt und $\mathcal{A}$ eine Familie von abgeschlossenen Mengen mit der endlichen Durchschnittseigenschaft.
	Angenommen $\bigcap_{A \in \mathcal{A}} A = \emptyset$.
	Dann gilt
	\[
		X = X \setminus \bigcap_{A \in \mathcal{A}} A = \bigcup_{A \in \mathcal{A}} X \setminus A.
	\]
	Nun ist $\mathcal{U} \coloneqq \set*{X \setminus A \given A \in \mathcal{A}}$ eine offene Überdeckung von $X$ und da $X$ kompakt ist, existiert $\mathcal{A}_0 \subset \mathcal{A}$ endlich, sodass
	\[
		X = \bigcup_{A \in \mathcal{A}_0} X \setminus A = X \setminus \underbrace{\bigcap_{A \in \mathcal{A}_0 } A }_{\neq \emptyset} \quad \light
	\]
	Für die umgekehrte Implikation sei nun $\mathcal{U} = \set{U_i}_{i \in I}$ eine offene Überdeckung von $X$.
	Angenommen für jede endliche Teilmenge $J \subseteq I$ gilt $X \neq \bigcup_{i \in J} U_i$.
	Betrachte nun $\mathcal{A} =  \set{X \setminus U_i}_{i \in I}$. Dann gilt nach Annahme
	\[
		\bigcap_{i \in J} X \setminus U_i = X \setminus \bigcup_{i \in J} U_i \neq \emptyset.
	\]
	Also hat $\mathcal{A}$ die endliche Durchschnittseigenschaft. Nach Vorraussetzung gilt dann
	\[
		\emptyset \not= \bigcap_{i \in I} X \setminus U_i = X \setminus \underbrace{\bigcup_{i \in I} U_i}_{= X} \quad \light \qedhere
	\]
\end{beweis}


\printindex
\printbibliography
\listoffigures
\todototoc
\listoftodos[To-do's und andere Baustellen]
\end{document}
